\chapter{Bounded source coefficients}
\label{ch:cert-2}

\section{Two quartic contacts}
\label{component:p4ec1a51262f9}
\label{p4ec1a51262f9:a-full-parent-case-for-a-finite-positive-d-commuting-edge}

When one measurement pair becomes commuting, a quadratic angular
coefficient can vanish. The next coefficient is then quartic. We treat
the two such contacts together: the source determines a critical centre,
and a finite polynomial expansion about that centre gives the required
Gram matrix. All coefficients in this section are bounded, $d$ is a
positive unit, and $\lambda=0$. The source is physical.
Write $(n,m,C,A,D)$ for the coefficients denoted
$(e_1,e_2,e_3,e_4,e_5)$ in the event order of SG1.

\paragraph{A finite commuting edge.}
Suppose the testing residues are $u=0$, $v=v_0\ne0$ and
$[W:X:Y:Z]=[1:0:0:0]$, and $m$ is infinitesimal.
On the residue source the determinant numerator is
\[
 Q(U,V)=UE(V)+U^2B(V),\quad
 E(V)=q_{10}+hV+q_{12}V^2,\quad
 B(V)=q_{20}+q_{21}V+q_{22}V^2.
\]
Assume $q_{10},q_{12},B(v_0^2)>0$ and $E(v_0^2)=0$.
The unit event amplitude $vW-X+uvY-uZ$ can be set to one,
so $W=(1+X-uvY+uZ)/v$. The opened event contributes the
constant $4m$. The state block at $u=0$ is
\[
 M_0(v)=\begin{pmatrix}
 2(d+2n+2\beta)(1+v^2)&0&0\\
 0&2d(1+v^2)+4(C+A)(1+v^2)+4D+4\alpha v^2&4v(D-\alpha)\\
 0&4v(D-\alpha)&4\alpha+4Dv^2
 \end{pmatrix}.
\]
Assume $M_0(v_0)\succ0$. In the applications below its lower
minor is positive because
$\alpha d+2\alpha C+2\alpha A+2\alpha D(1+v_0^2)
+Dv_0^2(d+2C+2A)>0$.

The substitution $X=u^2x$, $Y=ua$, $Z=ub$, $U=u^2$ defines
the polynomial $K$ by the exact identity
\[
 F(u^2x,ua,ub,u,v)=4m+u^2K(x,a,b,v,U).
 \tag{FE1}
\]
Its coefficients are Laurent polynomials in $v$. Substitution in
SG2 shows that $K_x=UE_x$ and
\[
 (E_x)_x\big|_{U=0}=4(d+2n+2\beta)(1+v^2)>0.
 \tag{FE2}
\]
The residue of $K$ at $U=0$ is independent of $x$.
Minimizing its positive $(a,b)$ block leaves
\[
 r(v)=\frac{16(1+v^2)^2E(v^2)}{v^2\det M_0(v)}.
\]
Since $E\ge0$ and has a zero at $v_0^2>0$,
$E(V)=q_{12}(V-v_0^2)^2$. Thus the Hessian in $(a,b,v)$
is positive definite at the residue zero. Equation FE2 then
determines $x$. The equations
$\mathcal E=(E_x,K_a,K_b,K_v)=0$ have an invertible
Jacobian in $q=(x,a,b,v)$. They define a source-field branch
$q_*(c,U)$. Its minimized profile $k(c,U)=K(q_*(c,U),U)$ satisfies
\[
 k(c_0,0)=0,\qquad
 \gamma=k_U(c_0,0)=
 \frac{16(1+v_0^2)^2B(v_0^2)}{v_0^2\det M_0(v_0)}>0.
 \tag{FE3}
\]
The last equality follows by taking the Schur complement of the
three-state block in the event-normalized parent. Its determinant
is $\det P/v^2$; the coefficients of $U$ and $U^2$ give
$r(v)$ and FE3 respectively.

\paragraph{A finite-ratio corner.}
\label{component:p03b551f2a9a0}
\label{p03b551f2a9a0:a-positive-d-case-at-a-finite-ratio-distance-four-corner}
Now let $\alpha$ be infinitesimal. At the residue source suppose
\[
 Q(U,V)=q_{20}U^2+hUV+q_{02}V^2+
 q_{21}U^2V+q_{12}UV^2+q_{22}U^2V^2,
\]
where $q_{20},q_{02}>0$. Choose a nonzero real $t_0$ such that
\[
 q_{20}t_0^4+ht_0^2+q_{02}=0,
 \qquad q_{21}t_0^4+q_{12}t_0^2>0.
\]
Normalize the unit amplitude $vY-Z$ to $-1$, giving $Z=1+vY$.
The substitution $X=vx$, $Y=vy$, $u=vt$, $W=w$, $V=v^2$
defines a polynomial $K$ by
\[
 F=4\alpha+v^2K(x,y,t,w,V).
 \tag{FR1}
\]
Its constant coefficient in $V$ is
\[
\begin{split}
 K_0={}&4nx^2+4m(w-x-t)^2+4C(tw-y)^2+4Ay^2\\
 &+4D(tw-y-1)^2+4\beta(x+t)^2\\
 &+2d\{(tw-y)^2+(x+t)^2-2tw\}+4\alpha t^2.
\end{split}
\]
Assume the residue quadratic form at $t_0$ has kernel
$(w_0,x_0,y_0,1)$ and a positive $(w,x,y)$ principal block
$M_J(t_0)$. Its minimized scalar is
$16(q_{20}t^4+ht^2+q_{02})/\det M_J(t)$.
Physicality makes the numerator $q_{20}(t^2-t_0^2)^2$.
The full Hessian in $q=(x,y,t,w)$ is therefore positive definite.
The source equations $\nabla_qK=0$ define $q_*(c,V)$, and
\[
 \gamma=\partial_VK(q_0,0)=
 \frac{16(q_{21}t_0^4+q_{12}t_0^2)}{\det M_J(t_0)}>0.
 \tag{FR2}
\]
This is again the next coefficient of the full determinant after
the positive state block is minimized. The testing chart has
$u,v$ infinitesimal and $[W:X:Y:Z]$ of residue $[w_0:0:0:1]$.
The actual testing ratio $u/v$ is unrestricted.

\paragraph{The common choice of centre.}
In either case write $S=U$ or $V$, and put $\kappa=k(c,0)$.
If $\kappa\ge0$, choose $S_c=0$ and $q_c=q_*(c,0)$.
If $\kappa<0$, solve
\[
 \mathcal E(q_c,S_c)=0,\qquad
 K(q_c,S_c)+S_cK_S(q_c,S_c)=0,
 \tag{FC1}
\]
where $\mathcal E=\nabla_qK$ in the corner case. The Jacobian
has the invertible normal block above and last pivot $2\gamma>0$.
The semialgebraic implicit-function theorem and real closed transfer
give a unique nearby source solution. The function
$k(c,S)+Sk_S(c,S)$ increases from the negative value $\kappa$,
so $S_c>0$. Realizing the corresponding source measurement and
state shows that
\[
 \mu=c_*+S_cK(q_c,S_c)\ge0,
 \qquad c_*=4m\ \hbox{or}\ 4\alpha.
 \tag{FC2}
\]
Every use of the implicit-function theorem here is within the
source field. The tensor identity uses only the resulting source
coordinates and finite polynomial divisions.

\paragraph{Full state expansion.}
At the edge set
\[
 N_s=(X-u^2x_c,Y-ua_c,Z-ub_c),\qquad N_v=u(v-v_c).
\]
Quadratic state dependence gives exactly
\[
 F=4m+UK(x_c,a_c,b_c,v,U)
       +UE_xN_x+uK_aN_y+uK_bN_z+N_s^TMN_s.
\]
The centre equations make each derivative vanish at $(v_c,U_c)$.
Difference division in $v$ and $U$ therefore gives only the
cross terms
\[
 N_xU(U-U_c),\quad N_yu(U-U_c),\quad N_zu(U-U_c),
 \quad N_xuN_v,\quad N_yN_v,\quad N_zN_v,
\]
and $N_v^2$ and $N_vu(U-U_c)$ from the scalar profile.
All coefficients are bounded. For negative powers of the unit
$v$, use $(v^{-j}-v_c^{-j})/(v-v_c)
=-\sum_{i=0}^{j-1}v^{-i-1}v_c^{-(j-i)}$.

At the corner set
\[
 N=(X-vx_c,Y-vy_c,u-vt_c,v(W-w_c)),\qquad z=W-w_c.
\]
The linear part is $v\nabla_qK(q_c,V)\cdot N$.
The rest belongs to $(N_1,N_2,N_3,vz)^2$: in SG2, a term
containing one $z$ and only one of the first three normals carries
$v$, and a term containing $z^2$ and no such normal carries $v^2$.
Terms with two of the first three normals already belong to this
ideal, with bounded $z$ left in their coefficients. This gives the
whole quadratic remainder without dividing by $v$.

In both cases the derivative coefficients are divisible by $S-S_c$.
Their cross coefficients with that last generator are infinitesimal.
The fixed-centre scalar profile has a double zero at $S_c$ when
$\kappa<0$. Finite division thus yields
\[
 \begin{array}{ll}
 \kappa\ge0:&F=c_*+\kappa S+G^THG,\quad G=(N,S),\\
 \kappa<0:&F=\mu+G^THG,\quad G=(N,S-S_c).
 \end{array}
 \tag{FC3}
\]
At the edge $N=(N_s,N_v)$ and the normal residue block is the
positive $X$ entry together with half the positive $(a,b,v)$
Hessian. At the corner it is half the positive four-variable
Hessian. The last diagonal is $\gamma>0$, and its cross entries
vanish in residue. The strict bounded-Gram lemma applies.
The remaining scalars in FC3 are nonnegative in the source field;
$S$ is a testing square. Restoring the event normalization completes
both certificates, including arbitrary relative infinitesimal rates.

For the edge with $E(v_0^2)>0$, the four-normal quadratic form in
$(X,Y,Z,u)$ is already strict, with $4m$ retained separately.
This supplies the other finite-edge alternative.


\section{\texorpdfstring{Positive total scale: the finite-pitchfork case}{Positive total scale: the finite-pitchfork case}}
\label{component:pfb0cf454f653}


The finite-edge calculation applies whenever $d,C,A,D,\alpha$
are positive units. We establish its hypotheses and then treat the
remaining commuting and projective limits.

\begin{theorem}
Let the original normalized C1 or C2 source be bounded, nonnegative and
physical over a real closed coefficient extension of $\mathbb R$.
Assume $d,C,A,D,\alpha$ are positive bounded units. Retain its exact
canonical zero coefficients, in particular $\lambda=0$, and every
other nonnegative original coefficient $n,m,\beta$. Then its entire
parent is tensor-SOS at every independent projective testing pair and
every real or complex testing state.
\end{theorem}

Noncommuting testing residues are covered by
Theorem~\ref{thm:noncommuting-refined}.

\subsection*{The finite edge}

The only possible zero at $u=0$, $v=v_0\ne0$ is the $W$ state
ray when $m$ is infinitesimal. We verify the conditions of
Section~\ref{component:p4ec1a51262f9} on that domain.

At $\lambda=m=0$, the exact determinant polynomial is
\[
 Q(U,V)=U E(V)+U^2 B(V),\quad
 E(V)=q_{10}+q_{11}V+q_{12}V^2,\quad
 B(V)=q_{20}+q_{21}V+q_{22}V^2.
\]
Use the positive quantities $t_1,\ldots,t_8$ in SG6.

All eight have positive bounded-unit residues. At the residue $m=0$,
the required coefficients are
\[
 q_{10}=4\alpha A t_7t_3>0,\quad
 q_{12}=4CDt_5t_2>0,\quad
 q_{21}=4(\beta n t_3t_2+CA t_1t_4)>0,
\]
and $q_{20},q_{22}\geq0$. Therefore $B(v_0^2)>0$ for every
$v_0\ne0$. Physicality implies $E(V)\geq0$ for all real $V>0$:
divide $Q(U,V)\geq0$ by $U>0$ and let the real variable $U$ tend to
zero. Consequently a zero $E(v_0^2)=0$ is a double zero,
$E(V)=q_{12}(V-v_0^2)^2$, and otherwise $E(v_0^2)>0$.

The state complement $M_0(v_0)$ in
Section~\ref{component:p4ec1a51262f9} is positive definite: its
lower determinant has the positive factor
$\alpha d+2\alpha C+2\alpha A+2\alpha D(1+v_0^2)
+Dv_0^2(d+2C+2A)$. Thus that section proves the full state
certificate at this edge, including both possible signs of the
optimized source value.



\subsection*{The other measurement limits}

Noncommuting residue tests use A6. It remains that at least one angle
has residue zero or infinity. For finite $u_0\ne0$ and $v_0=0$,
the real residue parent is a sum of nonnegative event terms and the
core $2d[(uW-Y)^2+(X+uZ)^2]$. The unit $A,\alpha,C$ terms
force $Y=Z=W=0$, and the core then forces $X=0$.
Hence the residue state matrix is positive definite.

For $u_0=0$ and finite $v_0\ne0$, the unit $C,A,D,\alpha$ terms
and $d>0$ leave only the $W$ ray if $m$ is weak, while unit $m$
removes that ray too. In the weak case, positive residue slack is
immediate; at the $W$ ray, the previous section applies for either
$E(v_0^2)>0$ or $E(v_0^2)=0$.

For $u=\infty$ and finite nonzero $v$, set $p=1/u$ and homogenize
as $F_p=p^2F(1/p,v)$. Its real residue has core
\[
 2d[(W+vX)^2+(vY-Z)^2]
\]
and positive $4C(1+v^2)W^2+4A(1+v^2)Y^2$ squares. These
force every state coordinate to vanish at a zero, so the matrix is
positive definite. For $v=\infty$ and finite nonzero $u$, the
homogenized residue has core $2d(1+u^2)(X^2+Y^2)$ together with
$4C(uW-Y)^2+4D(uX-Z)^2$. Again it is positive definite.
These are exact residue identities and use no density argument.

It remains to cover the four double corners. The following displayed
quadratic forms are complete, with coefficient letters interpreted in
real residue. All coefficients that are merely nonnegative remain so.
At a unit state coordinate, normalize it in the testing field.

At $(u,v)=(0,0)$ the only possible zero state is the $W$ ray.
Put $W=1$ and retain the entire $m$ square. The remainder lies
exactly in $(u,X,Y,Z)^2$, and its residue quadratic is
\[
 2(d+2n+2\beta)X^2+(4C+4D+2d)(u-Y)^2+4AY^2+4\alpha Z^2.
\]
It is strict. The testing $v$ is bounded infinitesimal in its
coefficients and is never divided out.

At $(u,v)=(\infty,0)$ the only possible zero state is the $X$ ray
when $n$ is weak; if $n$ is a unit there is no zero. Put $X=1$,
retain the entire $n$ square in $F_p$, and use roots $(p,v,W,Y,Z)$.
The whole remainder belongs to their square ideal and its residue is
\[
 4CW^2+4AY^2+(4D+2d)(W+v)^2
                 +4\alpha Z^2+[4(m+\beta)+2d](p+Z)^2\succ0.
\]

At $(u,v)=(0,\infty)$ the only possible zero state is the $W$ ray
when $m$ is weak. Put $q=1/v$, $W=1$, and retain the entire $m$
square in $F_q=q^2F(u,1/q)$. The whole remainder is in
$(u,X,Y,Z)^2$ and its residue quadratic is
\[
 (4n+4\beta+2d)X^2+4C(u-Y)^2+[4(A+\alpha)+2d]Y^2+4DZ^2\succ0.
\]
The testing $q$ remains a bounded infinitesimal coefficient.

At $(u,v)=(\infty,\infty)$ the only possible zero state is the $Z$
ray when $\beta$ is weak. Put $Z=1$ and retain the entire $\beta$
square in $F_{pq}=p^2q^2F(1/p,1/q)$. The complete remainder lies
in $(p,q,W,X,Y)^2$ and its residue quadratic is
\[
 4CW^2+4AY^2+4D(X-p)^2+(2d+4n)X^2
                            +[2d+4(m+\alpha)](Y-q)^2\succ0.
\]

For each corner, allocate two normal factors from every monomial to
the Gram roots; every remaining factor is bounded. Terms above
quadratic order have infinitesimal coefficients. The displayed full
quadratic residues are strict, so the separated strict-Gram lemma
applies. Restore each retained square. Homogeneous angle factors and
state normalizations are pure testing squares. The identities include
literal zero and infinity and every original nonnegative opening rate.
For nonzero state residues away from the listed rays, positivity gives
the scalar strict-residue exit. The zero state is immediate and complex
states split into real and imaginary evaluations. This proves the theorem.



\section{\texorpdfstring{Bounded C2 sources with infinitesimal alpha}{Bounded C2 sources with infinitesimal alpha}}
\label{component:pf31e38136e03}



\begin{theorem}
Let the original bounded normalized C2 source be physical, with
$\lambda=\beta=0$, all event coefficients nonnegative,
$d,n,m,C,A,D$ positive bounded units, and $\alpha$ weak, meaning
$\bar\alpha=0$. Then the complete original parent is tensor-SOS
for every independent projective testing pair and every real or
complex testing state.
\end{theorem}

The finite-ratio calculation of
Section~\ref{component:p03b551f2a9a0} treats the remaining corner.
We first identify the zero states to which it applies.

\subsection*{Only the zero-zero corner needs a new argument}

A6 covers noncommuting testing residues. On finite $u=0$, $v\ne0$,
the unit $m$ removes the $W$ ray, and the $(X,Y,Z)$ complement is
strict because $d,D,C,A$ are units and $v\ne0$. On finite $v=0$,
$u\ne0$, the units $n,A$ force $X=Y=0$ at a zero; the $C$ square
forces $W=0$, and the signed core $(X+uZ)^2$ forces $Z=0$.

The exact single-infinity core identities in
Section~\ref{component:pfb0cf454f653} require no unit $\alpha$: at
$u=\infty$, $v\ne0$ use
$2d[(W+vX)^2+(vY-Z)^2]$ and the $C,A$ squares; at $v=\infty$,
$u\ne0$ use $2d(1+u^2)(X^2+Y^2)$ and the $C,D$ squares.
Both state matrices are strictly positive.

At $(\infty,0)$ the unit $n$ removes the only possible $X$ zero
ray, and at $(0,\infty)$ the unit $m$ removes the $W$ ray.
At $(\infty,\infty)$, put $p=1/u,q=1/v$, normalize the possible
$Z$ zero ray to $Z=1$, and use $(p,q,W,X,Y)$ as roots. The exact
full-parent square-ideal allocation given in
Section~\ref{component:pfb0cf454f653} has residue
\[
 4CW^2+4AY^2+4D(X-p)^2+(2d+4n)X^2
                         +[2d+4(m+\alpha)](Y-q)^2\succ0.
\]
This form remains strict for weak $\alpha$. Thus every projective
test except $u,v$ infinitesimal is already covered, including literal
zero/infinite endpoints and all signs.

\subsection*{The exact two state charts at zero-zero}

Write $F$ for the original parent, with $\lambda=\beta=0$, and retain
the whole nonnegative $\alpha$ square
$4\alpha(1+u^2)(vY-Z)^2$ whenever using an ordinary strict jet.
At $u=v=0$ the residue is strictly positive on the $(X,Y)$ state
subspace. Hence $X$- or $Y$-unit states have positive residue slack.
For all other states, normalize a maximal coordinate. Then either
$W=1$ with $Z$ infinitesimal, or $Z=1$ with $W$ bounded.

In the first chart the entire remaining parent belongs to
$(X,Y,u,v)^2$. Its quadratic residue is
\begin{equation}
 (4n+2d)X^2+4m(v-X)^2+[4(C+D)+2d](u-Y)^2+4AY^2\succ0. \tag{1}
\end{equation}
Consequently the finite normal allocation and restoration of the
complete $\alpha$ square give tensor-SOS. No testing ratio is used.

In the second chart, set $w=\bar W\in\mathbb R$. The full residue
quadratic jet on $(X,Y,u,v)$ is
\begin{align}
 Q_w={}&4nX^2+4m(vw-X-u)^2+4C(uw-Y)^2+4AY^2\notag\\
      &+4D(uw-Y-v)^2
        +2d[(uw-Y)^2+(X+u)^2-2uvw].                    \tag{2}
\end{align}
All coefficient letters in (1)--(2) now denote their real residues.
Physicality implies $Q_w\succeq0$: evaluate the physical residue
parent at $(X,Y,u,v)=s(x,y,a,b)$ and let the real $s$ tend to zero.
The $(X,Y)$ principal block is strictly positive, already from the
$nX^2$ and $AY^2$ terms. If $Q_w$ is strict, the same ordinary
four-normal allocation applies, with $W-w$ infinitesimal in its
coefficients, and the whole $\alpha$ square restored.

Suppose instead that $Q_w$ is singular. It has no nonzero kernel
vector with $v=0$: the nonnegative sum at $v=0$ includes
$4nX^2+(4m+2d)(X+u)^2+4AY^2$, so $X=u=Y=0$.
It also has no nonzero kernel vector with $u=0$: the sum at $u=0$
includes $4nX^2+4AY^2+4D(Y+v)^2$, so $X=Y=v=0$.
Its real angular Schur complement therefore has positive diagonal
entries. Being singular and positive semidefinite, it has rank one
and a kernel with finite nonzero ratio $t_0=u/v$. Normalize a real
kernel vector by $v=1$ and write it as $(x_0,y_0,t_0,1)$.
This constructs a real coefficient-side jet center; it imposes no
condition on the actual independent testing value of $u/v$.

\subsection*{Application of the finite-ratio contact calculation}

The homogeneous state jet at a real ratio $t$ is
\begin{align*}
 J_t(W,x,y,Z)={}&4nx^2+4m(W-x-tZ)^2+4C(tW-y)^2+4Ay^2\\
 &+4D(tW-y-Z)^2
   +2d[(tW-y)^2+(x+tZ)^2-2tWZ].
\end{align*}
Its $(W,x,y)$ complement at $Z=0$ is strict: it contains
$4nx^2+4m(W-x)^2+4Ay^2$, with the remaining terms nonnegative.
At $t=t_0$, its kernel is $(w,x_0,y_0,1)$ by (2).

In the original determinant notation, at $\lambda=\beta=\alpha=0$,
\[
 q_{20}=4nAt_3t_4>0,\qquad q_{02}=4mDt_5t_6>0,
 \qquad q_{21}=4CAt_1t_4>0,\qquad q_{12}=4CDt_5t_2>0.
\]
Here the $t_i$ are the original positive sums of bounded coefficients
printed in the finite-pitchfork fit; all are units since $d$ is a
unit. The full jet determinant is
\[
 \det J_t=16(q_{20}t^4+q_{11}t^2+q_{02}).
\]
Thus the singular jet gives precisely the finite nonzero root required
by Section~\ref{component:p03b551f2a9a0}, and its next coefficient satisfies
$q_{21}t_0^4+q_{12}t_0^2>0$. The state complement, both endpoint
coefficients, and the positive next coefficient are all verified.

That calculation normalizes the pure-testing event
$vY-Z$ to $-1$, setting $Z=1+vY$. It keeps the complete opened
$\alpha$ square, including its $4\alpha u^2$ term, in the exact
profile
\[
 F=4\alpha+v^2K(x,y,t,w,V),\qquad
 (X,Y,u,W,V)=(vx,vy,vt,w,v^2)
\]
used only to construct coefficient-field critical centers. The positive
state complement and the double finite root give a strict full Hessian
in $(x,y,t,w)$, with positive next coefficient
\[
 \frac{16(q_{21}t_0^4+q_{12}t_0^2)}{\det M_J(t_0)}.
\]
The two signs of the optimized constant are handled by coefficient-only
implicit centers exactly as in the finite-edge proof; physicality makes
the remaining scalar nonnegative.

Its actual tensor roots are the un-divided expressions
\[
 X-vx_c,\quad Y-vy_c,\quad u-vt_c,\quad v(W-w_c),
 \quad v^2-V_c.
\]
All four full first-normal coefficients are $v\nabla K$ on this
frame. Monic difference division makes their coupling to $v^2-V_c$
infinitesimal. All 55 full-tail monomials have an exact allocation to
the first four roots; factors of $W-w_c$ that remain in coefficients
are infinitesimal. There is no substitution dividing $v(W-w_c)$ by
$v$. The full residue Gram is the positive four-variable Hessian and
the positive last scalar pivot. Hence the theorem applies to every
testing point with the stated residue state, regardless of whether its
actual tested $u/v$ is bounded or approaches $t_0$. This includes literal
$v=0$ and arbitrary independent opening rates.

The two zero-corner state charts and the strict/singular jet split
are exhaustive. Together with the complementary projective cases, this
proves the theorem. The zero state is direct, and complex states split
into real and imaginary quadratic evaluations.



\section{\texorpdfstring{The complete positive-scale C1 reciprocal case}{The complete positive-scale C1 reciprocal case}}
\label{component:p98f24864b071}

\label{p98f24864b071:sec:c1-reciprocal-rederived}

\paragraph{The source domain.}
Let the canonical C1 source be physical, with an actual finite
noncommuting zero contact, and normalize $d=1$. Suppose every source
coefficient is bounded, $\lambda=n=0$, $C,A,\alpha$ are positive
bounded units, $m>0$ and $D\geq0$ are infinitesimal, and
$\beta\geq0$ is arbitrary bounded. Retain the original canonical
source gaps and intersections. Then the entire parent has a finite
sum-of-squares representation in $K_c\otimes_{\mathbb R}K_t$ for
every independent projective test and every real or complex state.
The assertion includes literal $D=0$ and literal $\beta=0$ when
permitted by the actual contact; they cannot both vanish.

We use $O(1)$ and $o(1)$ for bounded and infinitesimal elements of
one of the two fields, and $a\asymp b$ for positive source quantities
whose ratio is a bounded unit. A bounded tensor Gram with a strictly
positive real residue is a sum of squares by the separated strict-Gram lemma. Its rows need not be bounded. No order on the tensor
product, mixed inverse, or cancellation of a determinant multiplier is
used below.

\subsection*{Source balances at the reciprocal corner}
\label{lem:reciprocal-source-balance}
Both C1 and C2 have a reciprocal corner when $m,D$ are
infinitesimal and $C,A,\alpha$ are positive units. All coefficients
are bounded, $d=1$, $\lambda=0$, and $n\beta=0$.
At a stationary source contact put $U=u_c^2$, $P=v_c^{-2}$.
The determinant formula SG5--SG6 becomes
\[
\begin{split}
 L(U,P^{-1})={}&h+\frac{ma_1}{U}+\frac{mDa_2}{UP}+a_3P\\
 &+\frac{\beta ma_4+Da_5}{P}+a_6U+na_8UP
 +\frac{\beta a_7U}{P},
\end{split}
\]
where the positive units are explicitly
\[
\begin{array}{llll}
a_1=4\alpha t_5t_8,&a_2=4t_5t_6,&a_3=4\alpha A t_7t_3,&a_4=4t_1t_6,\\
a_5=4Ct_5t_2,&a_6=4CAt_1t_4,&a_7=4Ct_1t_2,&a_8=4At_3t_4.
\end{array}
\]
Here $h$ is bounded. At a zero, every positive Laurent term is
bounded by $-h$, so $U,P$ are bounded. Differentiating at the
stationary contact gives the exact balances
\begin{equation}
\begin{aligned}
 U^2(a_6+na_8P+\beta a_7/P)&=m(a_1+Da_2/P),\\
 P^2(a_3+na_8U)&=mDa_2/U+\beta ma_4+Da_5+\beta a_7U.
\end{aligned}
\label{eq:reciprocal-source-balances}
\end{equation}
These equations force both $U$ and $P$ to be infinitesimal.
Indeed, if $U$ were a unit and $P$ a unit, the first balance
would have a positive-unit left side and an infinitesimal right
side. If $U$ were a unit and $P$ infinitesimal, the second balance
would give $mDa_2/P\le UP(a_3+na_8U)$, again making the first
right side infinitesimal. Thus $U$ is infinitesimal. If $P$ were
a unit, the first balance would then give
\[
 \frac{mDa_2}{U}
 =UP(a_6+na_8P)+\beta a_7U-ma_1P/U
 \le UP(a_6+na_8P)+\beta a_7U,
\]
which is infinitesimal. Every term on the second right side would
be infinitesimal, contradicting its positive-unit left side.

The nonnegative terms in the second balance now give
$D/P^2=O(1)$ and $\beta U/P^2=O(1)$. The first consequently gives
\begin{equation}
 \frac m{U^2}\asymp1+\frac\beta P,\qquad
 \frac mU,\ \frac DP,\ \frac{\beta U}P\in\operatorname{Inf}.
\label{eq:reciprocal-source-ratios}
\end{equation}
In C2, $\beta=0$, so $m/U^2$ is a positive unit and
\[
 \frac D{P^2}=\frac{a_3+na_8U}{a_5+ma_2/U}
\]
is a positive unit as well. Thus $m\asymp U^2$, $D\asymp P^2$.
For C1 the preceding bounded estimates also allow literal $D=0$
or $\beta=0$ whenever the source contact permits it.


\subsubsection*{Source entry before any state-ratio division}

Write the actual contact angles as $u_c=k>0$, $b_c=1/v_c=p>0$,
and put $U=k^2$, $P=p^2$. Equation~\eqref{eq:reciprocal-source-balances}, with $n=0$,
gives $U,P$ positive infinitesimals and
\[
 D/P^2=O(1),\qquad \beta U/P^2=O(1),\qquad
 m/U^2\asymp1+\beta/P.
 \tag{CR1}
\]
The same balances give
\[
 m/U=o(1),\qquad D/P=o(1),\qquad \beta U/P=o(1).
 \tag{CR2}
\]

At a physical noncommuting contact the signed core is negative, and
the exact core Schur identity SE9 gives $W_cZ_c>0$. Normalize $W_c=1$
and write $Z_c=q>0$. Define
\[
 Q=q^2,\quad L=PQ,\qquad
 r=\frac{k}{pq},\quad y=\frac{Y_c}{pq},\quad
 x=\frac{X_c}{pq^2}.
 \tag{CR3}
\]
We now prove these divisions are bounded wherever they are used.
In the reciprocal quarter-parent CR9 below the hard $X,Y$ matrix at
the source has residue
$\operatorname{diag}(\bar\beta+1/2,H_0)$, where
$H=C+A+\alpha+1/2$ is a positive unit. Eliminate this two-by-two
block inside the source field. Its Schur block on $W,Z$ satisfies
\[
 \frac{S_{WW}}{U}\sim a:=C-C^2/H>0,\qquad
 -\frac{S_{WZ}}{kp}\sim b_*:=1/2+C\alpha/H>0.
 \tag{CR4}
\]
Here and below a displayed positive expression in an asymptotic
formula is evaluated in residue. To check CR4 without a shape
assumption, the original hard couplings are
\[
\begin{aligned}
 P_{XW}&=p\{-m+(D+1/2)U\},&
 P_{YW}&=k\{m-C(1+P)-(D+1/2)P\},\\
 P_{XZ}&=k\{\beta-D+P(m+\beta+1/2)\},&
 P_{YZ}&=p\{D-\alpha-U(m+\alpha+1/2)\},\\
 P_{WZ}&=-(m+D+1/2)kp.
\end{aligned}
\]
CR2 shows that the $X$ corrections in CR4 vanish in residue and
the $Y$ corrections are respectively $C^2/H$ and $C\alpha/H$.
The kernel equation gives $q=S_{WW}/(-S_{WZ})$, so $r$ is a
positive bounded unit.

The scaled $X,Y$ kernel equations in the unknowns $x,y$ have
bounded entries and a diagonal positive residue. For example their
$X$ equation is exactly
\[
\begin{split}
 0={}&\{(\beta+1/2)(1+P)+Pm+(D+1/2)r^2PQ\}x\\
 &+r\{\beta-D+P[\beta+1/2-m(y-1)+(D+1/2)(r-y)]\}-m/Q.
\end{split}\tag{CR5}
\]
The quantity $m/Q=(m/U)r^2P$ is bounded before this equation is
used. The $y$ coefficient in CR5 is infinitesimal; the $x$
coefficient in the scaled $Y$ equation is $O(PQ)$ and its $y$
pivot is $H+o(1)$. Thus $x,y$ are bounded. Taking the $Y$ and
$Z$ kernel equations in residue, using CR1--CR2, yields
\[
 Hy\sim Cr+\alpha,\qquad C(r-y)\sim1/2,
 \qquad \alpha(1-y)\sim r/2.
 \tag{CR6}
\]
Consequently $y,r-y,1-y$ are positive bounded units. Equivalently
their real residues satisfy
\[
 C_0=\frac1{2(r_0-y_0)},\quad
 A_0=\frac{1+r_0-y_0}{2y_0},\quad
 \alpha_0=\frac{r_0}{2(1-y_0)}.
 \tag{CR7}
\]
For the last equation in CR6, divide the original $Z$ kernel
equation by $Pq$. Its extra $X$ contribution is
$O(x\beta Q+xPQ)$; CR1 and the already proved unit $r$ give
$\beta Q=O(P)$, while $PQ=U/r^2$ is infinitesimal. All the
omitted terms therefore vanish in residue. This proves CR6 in the
unbounded-$q$ case as well.

In CR5, $D=o(P)$ and $m/Q=o(P)$. Its numerator for $-x$ is
the sum of $r\beta$ and a positive-unit multiple of $P$, up to
$o(P)$. Hence the following bounds hold on the whole source domain:
\[
\begin{gathered}
 L=PQ=o(1),\qquad \beta Q/P=O(1),\qquad
 -x\asymp P+\beta,\qquad r+x\text{ is a positive unit},\\
 m\asymp PQ^2(P+\beta),\qquad
 m/(PQ)=O(PQ+\beta Q)=o(1),\qquad xQ=o(1).
\end{gathered}\tag{CR8}
\]
The statement about $r+x$ follows directly from CR5:
$x/r=-\beta/(\beta+1/2)+o(1)$.

\subsubsection*{The entire scalar and all normal coefficients}

Use $b=1/v$, $D_u=1+u^2$, $D_b=1+b^2$. The full reciprocal
quarter-parent is
\[
\begin{aligned}
 F={}&m(W-bX+uY-ubZ)^2+CD_b(uW-Y)^2+AD_uD_bY^2\\
 &+(D+1/2)(ubW+uX-bY-Z)^2
 +(\alpha+1/2)D_u(Y-bZ)^2\\
 &+(\beta+1/2)D_b(X+uZ)^2-D_uD_bZ^2/2.
\end{aligned}\tag{CR9}
\]
All further local identities are expansions of this full expression.
In the chart $W=1$, $Z=z$ put
\[
 u_g=rbz,\qquad X_g=xbz^2,\qquad Y_g=ybz,
 \qquad V=b^2,\quad T=z^2.
 \tag{CR10}
\]
Its entire scalar restriction is
\[
 \Phi=m+DT+hVT+\mathcal C V^2T+EVT^2+
 \mathcal F V^2T^2+GV^3T^2+JV^2T^3,
 \tag{CR11}
\]
where
\[
\begin{aligned}
 \mathcal C&=Ay^2+(C+D+1/2)(r-y)^2,\\
 E&=\beta(r+x)^2+x^2/2-2Drx,\qquad
 G=Ar^2y^2,\qquad J=(D+1/2)r^2x^2,\\
 \mathcal F&=Ar^2y^2+(\alpha+1/2)r^2(y-1)^2
       +\beta(r+x)^2+m[r(y-1)-x]^2\\
 &\hspace{12mm}+(2D+1)rx(r-y)+rx+x^2/2,\\
 h&=Ay^2+C(r-y)^2+\alpha(y-1)^2+y^2/2-r
          -2D(r-y)+2m[r(y-1)-x].
\end{aligned}\tag{CR12}
\]
In particular $\mathcal C,G$ are positive units,
$E\asymp\beta+P^2>0$, $J\geq0$, and all these coefficients
are bounded. If $\beta$ is infinitesimal, $\mathcal F$ is a
positive unit, with residue $r_0^2(1+r_0-y_0)/2$.
For unit $\beta$, $\mathcal F$ can have either sign; it is not
discarded or assumed positive.

Actual contact and its two angular/state jets give
$\Phi(P,Q)=\Phi_V(P,Q)=\Phi_T(P,Q)=0$. Therefore
\[
\begin{aligned}
 D&=P^2\mathcal C-PQE+P^3QG-P^2Q^2J,\\
 m&=PQ^2E+P^2Q^2\mathcal F+P^3Q^2G+2P^2Q^3J,\\
 h&=-2P\mathcal C-QE-2PQ\mathcal F-3P^2QG-2PQ^2J.
\end{aligned}\tag{CR13}
\]
Set
\[
 \delta=D/P^2\geq0,\qquad \theta=EQ/P>0.
 \quad \delta,\theta=O(1),\qquad
 \delta+\theta=\mathcal C+GL-JQ^2=\mathcal C+o(1),
 \tag{CR14}
\]
where $L=PQ$. The bounds follow either from CR1 and CR8 or from CR13;
$JQ^2=O((xQ)^2)=o(1)$. At least one of $\delta,\theta$
is a unit. Define
\[
 R=z(V-P),\quad K=VT-PQ,\quad
 M=V^2T+(P/Q)VT^2+P^2Q-3PVT.
\]
The exact complete balanced scalar is
\[
\begin{aligned}
 \Phi={}&\delta R^2+\theta M+\mathcal F K^2
       +G\{PK^2+VzRK\}\\
       &+J\{(QR-zK)^2+2QK^2\}.
\end{aligned}\tag{CR15}
\]
The entire $J$ line is retained as positive squares.

To retain the complete first-normal terms, Let $L_u,L_X,L_Y$
be half the first derivatives of CR9 at CR10. Then
\[
 L_u=bz A_u(V,T),\quad L_Y=bz A_y(V,T),\quad
 L_X=b\{-m+TB_X(V,T)\},
 \tag{CR16}
\]
where
\[
\begin{aligned}
 A_u&=a_{00}+a_{10}V+a_{01}T+a_{11}VT+a_{21}V^2T+a_{12}VT^2,\\
 A_y&=e_{00}+e_{10}V+e_{11}VT+e_{21}V^2T,\\
 B_X&=c_0+c_1V+c_2VT,\\
 c_0&=-Dr+\beta(r+x)+x/2,\\
 c_1&=(D+1/2)r(r-y)+(\beta+1/2)(r+x)
                        +m[x-r(y-1)],\\
 c_2&=(D+1/2)r^2x.
\end{aligned}\tag{CR17}
\]
Every displayed coefficient is bounded. The only bounds beyond this
needed below are
\[
 a_{01}=\beta(r+x)-Dx=O(E),\quad
 a_{12}=(D+1/2)rx^2,\quad a_{12}Q=O(1),
 \qquad c_2=O(P+E).
 \tag{CR18}
\]
They follow from CR8 and CR12. Every other coefficient is the corresponding monomial coefficient
of the explicit polynomial CR9.

State and angle stationarity give
$A_u(P,Q)=A_y(P,Q)=0$ and $B_X(P,Q)=m/Q$.
Put $E_T=bz(T-Q)$. Finite monic division therefore gives exactly
\[
\begin{aligned}
 L_i&=A_{i,V}(V,T)bR+B_i(T)E_T,\
 B_u(T)&=a_{01}+P\{a_{11}+Pa_{21}+a_{12}(T+Q)\},\\
 B_y(T)&=P(e_{11}+Pe_{21}),\\
 L_X&=b\{[m/(PQ)+c_2T]K+z[c_1-m/(PQ)]R\}.
\end{aligned}\tag{CR19}
\]
The first line is for $i=u,y$ and
$A_{i,V}=[A_i(V,T)-A_i(P,T)]/(V-P)$, an exact polynomial.
In particular, using $P(T-Q)=K-zR$, each of $L_u,L_Y,L_X$
is a bounded linear combination of
\[
 bR,\qquad bzK,\qquad b[m/(PQ)+c_2T]K,
 \qquad E E_T.
 \tag{CR20}
\]
Every occurrence of $K$ outside the last product has an
infinitesimal coefficient. CR19 remains literal at $b=0$ or $z=0$.

\subsubsection*{The cubic conductors and the bounded source-ratio chart}

All radicals in this subsection are products of radicals in the
separate fields. Set
\[
 \mathsf A=(V^2T)^{1/3},\quad
 \mathsf B=((P/Q)VT^2)^{1/3},\quad
 \mathsf C=(P^2Q)^{1/3},\qquad \rho=(P/Q)^{1/3}.
\]
Because $\mathsf A\mathsf B\mathsf C=PVT$,
\[
 M=\tfrac12(\mathsf A+\mathsf B+\mathsf C)
 [ (\mathsf A-\mathsf B)^2+(\mathsf A-\mathsf C)^2
                         +(\mathsf B-\mathsf C)^2].
 \tag{CR21}
\]
Let $\mathcal Z$ be its nine roots multiplied by $\sqrt\theta$,
so that $\theta M=\frac12\sum_{j=1}^9\mathcal Z_j^2$.
For nonnegative testing $b,z$, the following exact identities hold:
\[
\begin{aligned}
 bR={}&b\sqrt{\mathsf A}(\mathsf A-\mathsf B)
       +\rho b^{1/3}z^{2/3}\sqrt{\mathsf A}(\mathsf A-\mathsf C),\\
 E E_T={}&\sqrt E\,z\,[\sqrt\theta\sqrt{\mathsf B}
                                      (\mathsf B-\mathsf A)]\\
 &+\theta^{1/3}E^{1/6}b^{2/3}z^{1/3}
       [\sqrt\theta\sqrt{\mathsf B}(\mathsf B-\mathsf C)].
\end{aligned}\tag{CR22}
\]
Negative signs are handled by the exact testing gauges; absolute
values give the same monomial identities if desired. If $\theta$
is a unit, the first line has infinitesimal coefficients on
$\mathcal Z$: $\rho=(E/\theta)^{1/3}$ is bounded and $b$ is
infinitesimal. If $\delta$ is a unit, $bR$ instead has the
infinitesimal coefficient $b$ on the strict row $R$. The second
line has infinitesimal coefficients whenever $E$ is infinitesimal
or $z$ is infinitesimal. Its displayed coefficients use no
comparison of a source rate with a testing rate.

Suppose first that actual $q$ is bounded and $\beta$ is
infinitesimal. Use $K$ and $\mathcal Z$, with $R$ added when
$\delta$ is a unit. If only $\theta$ is a unit, retain
$\delta R^2$ as a separate positive square. In CR15 the $K$
pivot $\mathcal F$ is a positive unit. The only signed extra
term is
$GVzRK=Gbz(bR)K$, an infinitesimal Gram correction by CR22.
The term $GPK^2$ is nonnegative. Thus CR15 has a strict residue
Gram on the stated roots. CR19--CR22 give infinitesimal lifts of
all three complete first normals to these same roots.

If $\beta$ is a unit, CR8 forces $Q$ infinitesimal and $E$ a
unit. The only chart still requiring CR10 will be $z$ infinitesimal.
The second line of CR22 still has infinitesimal coefficients.
When $\delta$ is a unit, use the exact identity
\[
 2QM+TR^2-K^2=(K-zR)^2+2QR^2.
 \tag{CR23}
\]
Choose a positive real $a$ with $\overline{\mathcal F}+a>1$.
Subtracting $a$ times the corresponding zero identity changes the
first three terms of CR15 into
\[
 (\delta-aT+2aQ)R^2+(1-2aP/E)\theta M
       +(\mathcal F+a)K^2+a(K-zR)^2.
 \tag{CR24}
\]
The first three pivots have strictly positive residues; both costs
$T$ and $P/E$ are infinitesimal. This handles the signed
$\mathcal F$ exactly.

When instead $\theta$ is a unit, $\rho$ is also a unit and
\[
 K=\rho^{-1}\{\sqrt{\mathsf A}\,
       [\sqrt{\mathsf A}(\mathsf B-\mathsf C)]
       +\sqrt{\mathsf C}\,
       [\sqrt{\mathsf C}(\mathsf A-\mathsf C)]\}.
 \tag{CR25}
\]
Both coefficients on $\mathcal Z$ are infinitesimal. Thus the
entire signed $\mathcal F K^2$ is an infinitesimal correction on
the strict circuit Gram. The $G$ cross term and every first normal
are controlled by CR22 and CR25; $\delta R^2$ and the $J$ line
remain positive squares. No inverse of $\delta$ is used.

\subsubsection*{An unbounded actual state ratio in the same testing chart}

Now let actual $q$ be unbounded. CR8 and CR12 imply
\[
 \beta/P=O(1/Q)=o(1),\quad -x\asymp P,\quad
 m\asymp L^2,\quad E/P=o(1),\quad JQ^2=O(L^2)=o(1).
 \tag{CR26}
\]
Here $\beta$ is infinitesimal and $\mathcal F$ is a positive
unit. In the testing chart $W=1$, $z$ bounded, the exact relation
\[
 \theta M=\theta R^2+(E/P)(K^2-zRK)
 \tag{CR27}
\]
makes CR15 a strict two-row Gram on $R,K$, plus its whole positive
$J$ line. Its residue is
$\operatorname{diag}(\overline{\delta+\theta},\overline{\mathcal F})$.
Every correction coefficient is infinitesimal: $E/P$, $P$ and
$Vz$ are infinitesimal. In CR19, $B_u/P$ and $B_y/P$ are bounded,
because $a_{01}/P=O(\beta/P+Dx/P)$ and $a_{12}Q$ are bounded.
Using $P E_T=bzK-bz^2R$ shows that all three full normals are
infinitesimal linear combinations of $R,K$.

\subsubsection*{The complementary testing state chart}

Use $Z=1$, $W=w$ infinitesimal, and write $a=u/r$. The graph is
\[
 b=aw,\qquad X=xa,\qquad Y=yaw.
 \tag{CR28}
\]
The entire scalar is
\[
 \Phi'=D+mw^2+Ea^2+ha^2w^2+\mathcal C a^4w^4
       +\mathcal F a^4w^2+Ga^6w^4+Ja^4.
 \tag{CR29}
\]
First assume actual $q$ bounded and $\beta$ infinitesimal. If
$\delta$ is a unit, use the three strict rows
$P,a^2w,\sqrt E\,a$. The sole signed cross in CR29 is
$h a^2w^2=(h/P)w\,P(a^2w)$, and $h/P$ is bounded by CR13.
Every other term is positive or has the positive-unit $\mathcal F$
pivot. If $\theta$ is a unit, $P/E=Q/\theta$ is bounded;
use instead $\sqrt E\,a,a^2w$ and retain $D,mw^2$ whole.
The same signed term is the infinitesimal correction
$(h/E)w^2(\sqrt E\,a)^2$.

Here are the complete complementary first normals, which also prove
the bounds for unbounded actual $q$. Define
\[
 R'=a^2w^2-P,\qquad K'=w(a^2-L).
 \tag{CR30}
\]
Half the $b,Y$ derivatives of the original full CR9 at CR28 have
the respective forms
\[
 L_i'=aw(d_{i0}+d_{i1}a^2+d_{i2}a^2w^2+d_{i3}a^4w^2),
 \qquad i=b,Y,
 \tag{CR31}
\]
with bounded coefficients. At the actual contact
$a^2=L,w^2=1/Q$, they satisfy
$d_{i0}+Ld_{i1}+Pd_{i2}+PLd_{i3}=0$.
Consequently the entire, exact divisions are
\[
 L_i'=(d_{i1}+Pd_{i3})aK'
                   +(d_{i2}+a^2d_{i3})awR'.
 \tag{CR32}
\]
The full $X$ derivative is
\[
 L_X'=a(c_0+c_2a^2-mw^2+c_1a^2w^2),\qquad
 c_0+c_2L-m/Q+c_1P=0.
 \tag{CR33}
\]
For bounded $Q$, $c_0/P$ and $m/P$ are bounded, and
$c_2\in P B_c+E B_c$ by CR18. The products $aK'$ and
$awR'$ are infinitesimal linear combinations of $P,a^2w$.
Expanding CR33 gives infinitesimal coefficients on
$P,a^2w,\sqrt E\,a$. If $\theta$ is a unit, use $P/E=O(1)$:
$c_0/\sqrt E=o(1)$, $c_2=O(E)$ and $m/E=O(1)$, giving the
same assertion on the two rows $\sqrt E\,a,a^2w$.
These are pure-source bounded quotients; no testing $w$ is inverted.

For unbounded actual $q$, CR29 has the exact alternative form
\[
\begin{aligned}
 \Phi'={}&\delta R'^2+\theta M'+\mathcal F K'^2
       +G\{PK'^2+a^2wR'K'\}\\
       &+J\{a^4(Qw^2-1)^2+2QK'^2\},\\
 M'={}&a^4w^4+(P/Q)a^2+P^2Qw^2-3Pa^2w^2.
\end{aligned}\tag{CR34}
\]
Define nonnegative elementary radicals
\[
 \mathsf A'=|aw|^{4/3},\quad
 \mathsf B'=(P/Q)^{1/3}|a|^{2/3},\quad
 \mathsf C'=(P^2Q)^{1/3}|w|^{2/3}.
\]
They give the same nine-square cubic identity for $M'$. The exact
conductor needed for the complete normals is
\[
 |a|R'=|a|\sqrt{\mathsf A'}(\mathsf A'-\mathsf C')
          +\sqrt L\sqrt{\mathsf B'}(\mathsf A'-\mathsf B').
 \tag{CR35}
\]
Its coefficients are infinitesimal. If $\delta$ is a unit, use
$R',K'$ and retain the circuit as positive squares. If $\theta$
is a unit, use its scaled nine roots and $K'$, retaining
$\delta R'^2$ whole. In either case CR34 is strict: its sole
signed cross is $Gaw(aR')K'$, controlled by CR35, and its last
line is retained as squares.

CR26 implies $c_2/P=O(1)$ and
$c_0/P=m/(PQ)-Qc_2-c_1=O(1)$. The full exceptional normal
has the exact division
\[
\begin{aligned}
 L_X'={}&-[c_0/P+(c_2/P)a^2]\,aR'\\
 &+[(c_2/P)(a^2+L)+c_0/P+c_1]\,awK'.
\end{aligned}\tag{CR36}
\]
Together with CR32, this uses only bounded multiples of
$aR',awR',aK'$. Each has infinitesimal coefficients on the strict
soft rows by CR35. This also proves the result at literal $w=0$.

\subsubsection*{All higher normals and the remaining unit-beta states}

The preceding scalar and first-normal formulas certify the entire
parent. Here are the required hard blocks and the justification for
every higher term. Normalize the testing state so its coordinates
are bounded with at least one unit. At the reciprocal angle corner
$u,b=o(1)$ the state residue hard block is
$(\bar\beta+1/2)X^2+H_0Y^2$. If $X$ or $Y$ is a unit, the
entire parent has positive real residue and the strict-element lemma
applies immediately. Otherwise $X,Y$ are infinitesimal and one of
$W,Z$ is a unit.

The exact sign gauges make $u,b\geq0$. If $WZ\leq0$, the signed
core is a nonnegative pure-testing scalar by the full core Schur
identity SE9 and its literal homogeneous endpoints, so its source
multiple and all event squares are separated squares. Otherwise
choose either $W=1$, $z=Z$ bounded, or $Z=1$, $w=W$
infinitesimal. Both have $X,Y$ infinitesimal.

In the $W=1$ chart, set
$N_u=u-rbz$, $N_X=X-xbz^2$, $N_Y=Y-ybz$.
They are infinitesimal. The complete expansion of CR9 is its scalar
CR11, twice the three first normals CR16, and a polynomial quadratic
Gram in $N_u,N_X,N_Y$. Terms of normal degree at least three
have infinitesimal Gram entries. For weak $\beta$, or for unit
$\beta$ with $z$ infinitesimal, the hard residue is
\[
 C_0(N_u-N_Y)^2+(A_0+\alpha_0+1/2)N_Y^2
                         +(\beta_0+1/2)N_X^2\succ0.
 \tag{CR37}
\]
All first-normal lifts established above have infinitesimal
coefficients. The full hard-and-soft residue is therefore positive
definite. In the $Z=1$ chart use
$N_b=b-aw$, $N_X=X-xa$, $N_Y=Y-yaw$.
For weak $\beta$ its hard residue is
\[
 N_X^2/2+H_0N_Y^2-2\alpha_0N_bN_Y+\alpha_0N_b^2\succ0,
 \tag{CR38}
\]
since $H_0-\alpha_0=C_0+A_0+1/2>0$.
The same finite expansion handles every higher normal term. Thus
none of CR15, CR27 or CR34 is being used merely as a scalar
restriction.

It remains only unit $\beta$ and a testing state with $Z$ a unit
relative to $W$. Normalize $Z=1$ and $W=\omega$ bounded. Retain
the entire $m,D$ squares. The remainder belongs exactly to the
square of the ideal $(u,b,X,Y)$. Its quadratic real residue is
\[
\begin{aligned}
 &(\beta_0+1/2)\left(X+\frac{\beta_0}{\beta_0+1/2}u\right)^2
      +\frac{\beta_0}{1+2\beta_0}u^2\\
 &\quad+H_0\left(Y-\frac{C_0\omega u+\alpha_0 b}{H_0}\right)^2
      +a_0(\omega u-r_0b)^2,
 \qquad a_0=C_0-C_0^2/H_0>0.
\end{aligned}\tag{CR39}
\]
CR7 is precisely the critical relation needed for this identity.
The expression is positive definite in $(u,b,X,Y)$ for every
real $\omega$, including zero. Its complete bounded polynomial
matrix is consequently strict; every higher degree is an
infinitesimal correction. This completes all local states, including
the unit-beta case where $\mathcal F$ was signed.

\subsubsection*{Every other test and restoration}

The exact outer identities RC18--RC25 from the newly rederived
C2 reciprocal theorem apply here as polynomial identities, with
$n=0$ and the entire extra $\beta$ square retained. Their strictness
uses only $C,A,\alpha$ positive units and the critical residue
relations CR7; no C2 source membership is used. Explicitly:
\begin{itemize}
\item For bounded original $u,v$, the four rows $(uW,Y,X,Z)$
give RC18. After CR7 its real determinant is
$ (r_0+u^2v^2)D_u^2D_v^2(1+r_0-y_0)^2/
[16y_0(r_0-y_0)(1-y_0)]>0$.
\item For Alice reciprocal and Bob bounded, use the four rows
$(W,Y,vX,sX+Z)$ and the full core RC20. Its real residue RC21
is strict using $C,A>0$.
\item For Bob reciprocal and Alice a bounded unit, use
$(uW-Y,Y,X,bZ,\sqrt D\,Z)$ and the entire $D$ square as in
RC22--RC23. CR7 makes the three-by-three angular determinant
strict. Literal $D=0$ causes its last row to be zero, which is
allowed in a strict coefficient Gram.
\item For both angles reciprocal, RC24--RC25 give the five rows
$(W,X,Y,bZ,\sqrt D\,sZ)$ and a strictly positive real residue.
\end{itemize}
The only remaining angle chart is exactly $u,b=o(1)$, treated
above. These are whole-parent identities with every omitted event
square retained, and include all literal axes and infinities.
The exact source and testing sign gauges, positive pure-field
homogeneous normalization factors and real/imaginary state splitting
restore the stated projective tensor-SOS theorem.



\section{Exhaustion of the bounded source domain}
\label{component:pc0b2849f8d09}

\label{pc0b2849f8d09:sec:positive-scale-union-rederived}

\paragraph{Exact statement.}
Let \(K_c,K_t\) be labelled real closed extensions of \(\mathbb R\).
The original canonical coefficient vector is physical over \(K_c\),
has nonnegative coefficients, and has an actual finite noncommuting
zero contact. In C1 the exact omissions are \(\lambda=n=0\),
with \(m,C,A,\alpha>0\), \(D,\beta\geq0\), \(D+\beta>0\);
in C2 they are \(\lambda=\beta=0\), with
\(m,C,A,D,\alpha>0\), \(n\geq0\).
Use the actual native sign chart of this contact, with its explicitly
positive source gaps, as supplied by canonical entry.
Assume that normalization by \(d>0\) makes all coefficients bounded.
Equivalently the normalized total-scale coordinate
\(d/(d+\lambda+n+m+C+A+D+\alpha+\beta)\) is a positive unit.
Then the entire original homogeneous slack is a sum of squares
in \(K_c\otimes_{\mathbb R}K_t\) for every independent projective
testing pair and every real or complex testing state.
Chapter~\ref{ch:entry} supplies these contact coordinates for every
source in the two families.

Division by \(d\) is a positive pure-source normalization.
Hence take \(d=1\) below. Every nonnegative bounded source quantity
is either infinitesimal (including literal zero where permitted)
or a positive unit. These two alternatives give the source partition below.

We use the following two cases established above:
\[
\begin{array}{c|l|l}
 &\text{positive-unit coefficients}&\text{proof}\\\hline
\mathrm S&d,C,A,D,\alpha&\text{Section~\ref{component:pfb0cf454f653}}\\
\mathrm{CU}&d,n,m,C,A,D\quad(\alpha\text{ infinitesimal})
 &\text{Section~\ref{component:pf31e38136e03}}.
\end{array}
\]
Every other nonnegative coefficient may have arbitrary residue in
the stated canonical support. Both proofs certify the full state
polynomial at all independent tests. Their noncommuting case is
Theorem~\ref{thm:noncommuting-refined}; their finite-edge and
finite-ratio cases were proved at the start of this chapter.


\subsection*{Symmetries and the source partition}

Both transformations below preserve the core and nonnegative event
coefficients, actual physicality, finite nonzero contact, and pure
total-scale normalization. They act separately in source and testing
fields, so transport a full tensor certificate.

For C1 let
\[
 \sigma_1=(\lambda\ n)(m\ C)(A\ \alpha)(D\ \beta),\qquad
 G_v(W,X,Y,Z)=(vW-X,-W-vX,Z-vY,Y+vZ).
\]
The entire raw parent has the polynomial identity
\[
 u^2F_{\sigma_1c}(1/u,v;G_v\psi)
       =(1+v^2)F_c(u,v;\psi),\qquad
 G_v^{\mathsf T}G_v=(1+v^2)I .
 \tag{AU12}
\]
Homogenization interchanges the two Alice angle coordinates;
the Bob-homogeneous \(G\) has norm factor \(v_0^2+v_1^2>0\).
Thus literal endpoints are transported without dividing zero.

For C2 let
\[
 \sigma_2=(\lambda\ \beta)(m\ D)(C\ \alpha),\qquad
 J(W,X,Y,Z)=(Z,X,Y,W).
\]
The full identity is
\[
 u^2v^2F_{\sigma_2c}(1/v,1/u;J\psi)=F_c(u,v;\psi).
 \tag{AU13}
\]
It is the physical party interchange with both first observables
negated, followed by the local real \(X\) basis changes.
The homogeneous identity includes both projective endpoints.
In particular the source squared half-angles transform as
\((U,V)\mapsto(1/V,1/U)\).

The two actual forbidden weak faces are
\[
\begin{aligned}
 h|_{\lambda=\beta=A=0}
 &=-(1+2n)(2D+2C+1)(2\alpha+1+2m),\\
 h|_{\lambda=n=m=C=0}
 &=-(2D+1)(2\beta+1)(2A+2\alpha+1).
\end{aligned}\tag{AU14}
\]
In C2 with \(A\) weak, the only positive Laurent terms are
\(q_{01}/U+q_{02}V/U+q_{12}V\), which tend to zero along
\(U=t^2,V=1/t,\ t\to\infty\).
In C1 with both \(m,C\) weak the only positive term is \(q_{10}/V\),
which tends to zero as \(V\to\infty\).
The negative constants contradict physicality.
Thus C2 has \(A\) unit, and C1 has at least one of \(m,C\) unit.


\paragraph{C2.}
If \(\alpha\) is weak, WA1--WA6 force \(C,A,D\) to be units.
There are precisely three disjoint alternatives:
\[
\begin{array}{c|c}
 n,m\text{ units}&\mathrm{CU}\\
 n\text{ unit},\ m\text{ weak}&\mathrm{WA1\!-\!WA27}\\
 n\text{ weak},\ m\text{ arbitrary}&\mathrm{WN1\!-\!WN44}.
\end{array}\tag{AU15}
\]
If \(C\) is weak, apply \(\sigma_2\) and the same table.
Every unit requirement in the transformed table is a consequence
of the transformed physical source, including its required
\(D'=m\) bound; it is not assumed for an excluded original case.

It remains that \(C,A,\alpha\) are units.
The actual dictionary then has
\[
 q_{01}=ma_1,\ q_{02}=mDa_2,\ q_{10}=a_3,\
 q_{12}=Da_4,\ q_{20}=na_5,\ q_{21}=a_6 ,
\]
where every \(a_i\) is a positive bounded unit.
The contact equation bounds \(U,V^{-1},m/U,DV\).
The two angular equations give exactly
\[
 \begin{split}
 DV^2(a_4+a_2m/U)&=a_3+na_5U,\\
 U^2(a_6+na_5/V)&=m(a_1+a_2DV).
 \end{split}
\]
All parentheses are positive bounded units, so
\(DV^2\asymp1\), \(U^2\asymp m\).
Consequently the exhaustive partition is
\[
\begin{array}{c|c}
 D\text{ unit},\ m\text{ arbitrary}&\mathrm S\\
 D\text{ weak},\ m\text{ unit}&\sigma_2(\mathrm S)\\
 D,m\text{ weak}&\mathrm{RC1\!-\!RC25}.
\end{array}\tag{AU16}
\]
RC is the direct full-parent C2 reciprocal theorem; it permits
arbitrary \(n\geq0\). No unit-\(n\) condition has entered this
part of the partition.

\paragraph{C1.}
If \(\alpha\) is weak, WN applies directly, with all its common
source-entry implications proved in WA1--WA10.
If \(A\) is weak, apply \(\sigma_1\), making \(\alpha'=A\) weak,
and use WN on the transformed physical source.
Otherwise \(A,\alpha\) are units.
By AU14, orient with \(\sigma_1\), if necessary, so that \(C\)
is a unit. The exact dictionary becomes
\[
\begin{gathered}
 q_{01}=ma_1,\ q_{02}=mDa_2,\ q_{10}=a_3,\quad
 q_{12}=\beta ma_4+Da_5,\quad q_{21}=a_6,\quad q_{22}=\beta a_7,\\
 U^2(a_6+\beta a_7V)=m(a_1+a_2DV),\\
 a_3=V^2\{mDa_2/U+\beta ma_4+Da_5+\beta a_7U\}.
\end{gathered}\tag{AU17}
\]
The contact equation again bounds \(U,V^{-1},m/U,DV\).
If \(m\) is weak, \(U\) is infinitesimal.
If \(m\) is a unit, \(U\) is a positive unit.
Indeed, the contact equation also bounds \(\beta UV\); if
\(U\) were infinitesimal, then
\(U^2(a_6+\beta a_7V)\) would be infinitesimal, contradicting
the positive unit \(ma_1\) on the right of AU17.
For unit \(m\), the last braces are comparable with \(D+\beta\).
For weak \(m\), bounded \(V\) forces \(m/U\) weak in the
preceding equation, and then the last equation forces \(D\)
unit; conversely unit \(D\) forces bounded \(V\).
Thus the following source partition is exhaustive:
\[
\begin{array}{c|c}
 D\text{ unit},\ m\text{ arbitrary}&\mathrm S\\
 D\text{ weak},\ m\text{ unit},\ \beta\text{ unit}
       &\sigma_1(\mathrm S)\\
 m\text{ unit},\ D,\beta\text{ weak}
       &\mathrm{PP1\!-\!PP19}\\
 m,D\text{ weak},\ \beta\text{ arbitrary}
       &\mathrm{CR1\!-\!CR39}.
\end{array}\tag{AU18}
\]
In the second line all of \(m,C,A,\alpha\) are units before
reflection, so the reflected strong theorem has every required
unit. PP is the C1 projective-pitchfork theorem, with actual
contact \(U\) unit and \(V\) unbounded.
CR is the C1 reciprocal all-tests theorem, including unit
\(\beta\) and literal allowed \(D=0\) or \(\beta=0\).
The last line treats infinitesimal \(m,D\).
No boundary is promoted by a density argument.

AU15--AU18 cover every permitted bounded positive-\(d\) canonical
source. Each leaf proves the entire original parent at every
independent projective test, with the exact same source coefficients.
The full-parent symmetries, positive pure-source \(d\) scaling,
positive pure-testing homogeneous factors and complex-state split
restore the assertion in the opening paragraph. This completes
A5 at its canonical-entry interface.

The source estimates in each listed case use precisely its
physical contact equations and the displayed coefficient bounds.
The finite partition therefore proves the all-tests assertion for
every bounded canonical source.

\section{A three-event face and its stationary zeros}
\label{component:p23a094ff816d}

\label{p23a094ff816d:sec:rederived-three-event-contact}

The three-event face has an explicit matrix representation. Its
dual gives a separated sum of squares. We then use the stationary
equations to show that this face contains no source with a finite
noncommuting contact.

\subsubsection*{The selected full parent}

Write $D_u=1+u^2$, $D_v=1+v^2$, and retain the coefficient order
$(d,\lambda,n,m,C,A,D,\alpha,\beta)$.  The original complete parent is
\begin{align}
F={}&4\lambda D_u(W+vX)^2+4nD_uD_vX^2
 +4m(vW-X+uvY-uZ)^2\notag\\
&+4CD_v(uW-Y)^2+4AD_uD_vY^2
 +(4D+2d)(uW+uvX-Y-vZ)^2\notag\\
&+(4\alpha+2d)D_u(vY-Z)^2
 +(4\beta+2d)D_v(X+uZ)^2-2dD_uD_vZ^2.
\tag{CAP1}
\end{align}
All nine coefficients are nonnegative.  Let $F_*$ be the selected parent
with $\lambda=n=m=\alpha=\beta=0$, so that only the consecutive event
triple $(D,C,A)$ remains.  For $d,C,A,D\geq0$ define
\[
\Delta=8ADC-4ADd-d^2\{d+2(C+A+D)\}.
\tag{CAP2}
\]

\begin{proposition}[Exact physical wall and strict interior-angle positivity]
If $d>0$, the selected parent $F_*$ is positive semidefinite at every
projective measurement pair if and only if $\Delta\geq0$.  Whenever
this holds, its state matrix is positive definite at every finite pair
$u,v\ne0$.  In particular, the equality wall $\Delta=0$ has no finite
noncommuting contact.
\end{proposition}

\begin{proof}
Let $P_*(u,v)$ be the symmetric state matrix.  Direct expansion of the
full polynomial (CAP1), without a state or angle truncation, gives
\[
\det P_*=16D_u^2D_v^2u^2v^2
 \{d\Delta+4CDd(d+2A)v^2+4CAd(d+2D)u^2\}.
\tag{CAP3}
\]
Its restriction to $Z=0$ is the sum of the
nonnegative event squares and
\[
2d\{(uW+uvX-Y)^2+D_uv^2Y^2+D_vX^2\}.
\tag{CAP4}
\]
For $uv\ne0$ these three squared linear forms are independent.  Hence
the upper $3\times3$ state block is positive definite.  A Schur
complement taken in this one ordered field proves that $P_*\succeq0$
if and only if its determinant is nonnegative, and that a positive
determinant gives positive definiteness.  This is a statement about
ordinary field-valued source or testing matrices; it takes no inverse
in a tensor product.

Suppose $\Delta<0$.  Put $r=4CDd(d+2A)$ and
$s=4CAd(d+2D)$, both nonnegative.  Taking
\[
u^2=v^2=\frac{-d\Delta}{2(r+s+1)}>0
\]
makes the expression in braces in (CAP3) negative.  This is an explicit
ordered-field witness to nonphysicality, so no limiting argument is
needed.  Conversely, $\Delta\geq0$ makes all three terms in braces
nonnegative.  It also forces $C,A,D>0$: if one vanished, (CAP2) would
be strictly negative because $d>0$.  Thus $r,s>0$ and the determinant
is strictly positive whenever $uv\ne0$.

At an actual commuting angle, including a projective infinity, one
party's two reflections commute.  Simultaneous diagonalization makes
the CHSH operator on each local eigenspace either $\pm2$ times one
remote reflection.  The signed core $I-S/2$ is then positive
semidefinite.  Every event square is positive semidefinite too.  This
covers all literal endpoints directly.  The argument holds over every
real closed field, because it uses real diagonalization of symmetric
matrices and explicit ordered-field identities.  If $d=0$, every
parent in (CAP1) is its event-square sum separately.
\end{proof}

\subsubsection*{A small exact lift of the selected source cone}

Set
\[
p=4C+4D+2d,\qquad r=4D+2d,\qquad
B(d,C,A,D)=
\begin{pmatrix}
p+4A&4D&-p\\
4D&4D&-r\\
-p&-r&p
\end{pmatrix}.
\tag{CAP5}
\]
For nonnegative coefficients the selected source is physical exactly
when $B\succeq0$.  Here is a direct verification of this matrix
description.  When $d>0$, use the positive last diagonal entry $p$.
The Schur complement of that entry is
\[
\begin{pmatrix}4A&-2d\\-2d&4b/p\end{pmatrix},
\qquad b=4CD-2Dd-d^2,
\qquad \det B=8\Delta.
\tag{CAP6}
\]
Moreover
\[
\Delta=2Ab-d^2(d+2C+2D).
\tag{CAP7}
\]
If $\Delta\geq0$, then $A,b>0$, and (CAP6) is positive semidefinite.
If $B\succeq0$, its determinant is nonnegative, hence $\Delta\geq0$.
For $d=0$ one has the exact quadratic identity
\[
(y,z,x)B(0,C,A,D)(y,z,x)^T
 =4C(y-x)^2+4Ay^2+4D(y+z-x)^2.
\tag{CAP8}
\]
Thus the entire selected physical cone, including all zero faces, is
\[
\mathcal C=\{(d,C,A,D)\in\mathbb R_{\geq0}^4:B(d,C,A,D)\succeq0\}.
\tag{CAP9}
\]
It has a strict feasible point $(1,2,2,2)$: the leading principal
minors of $B$ there are $26,144,280$.

\begin{lemma}[Separated certificate for the selected wall]
For arbitrary labelled real closed fields $K_c,K_t\supset\mathbb R$,
every source $c\in\mathcal C_{K_c}$, and every projective angle pair
and real state in $K_t$, the complete selected parent is a finite sum
of squares in $K_c\otimes_{\mathbb R}K_t$.  The full parent (CAP1)
has the same property if its selected coefficients satisfy (CAP9).
\end{lemma}

\begin{proof}
We first establish the exact dual representation used in this lemma.
For a real linear functional $\ell$ nonnegative on $\mathcal C$,
there are $\mu\in\mathbb R_{\geq0}^4$ and a real symmetric
$Q\succeq0$ such that, for all $c=(d,C,A,D)$,
\[
\ell(c)=\mu\cdot c+\operatorname{tr}(B(c)Q).
\tag{CAP10}
\]
Indeed let $\mathcal D$ be the cone of functionals on the right.  It
is closed: if a sequence of these functionals converges, evaluation
at $c_0=(1,2,2,2)$ bounds both $\sum\mu_i$ and
$\operatorname{tr}Q$, because all coordinates of $c_0$ and the
smallest eigenvalue of $B(c_0)$ are strictly positive.  Positive
semidefinite matrices with bounded trace are bounded, so there is a
convergent subsequence of the representing pairs.  The dual cone
$\mathcal D^*$ is exactly (CAP9).  Finite-dimensional separation of
closed convex cones gives $\mathcal D=\mathcal C^*$ and proves
(CAP10).  This statement has a fixed finite first-order formulation
using the principal-minor inequalities for a $3\times3$ matrix.
Consequently it holds over every real closed field by transfer.

Fix an arbitrary testing tuple in $K_t$.  Its homogeneous selected
parent is a linear functional $\ell_t$ on the four source
coefficients.  The physical equivalence just proved says that
$\ell_t\geq0$ on $\mathcal C_{K_t}$.  Apply (CAP10) in $K_t$ to
obtain $\mu_t\geq0$, $Q_t\succeq0$.  Evaluating the resulting
linear coefficient identity in separated fields gives
\[
F_*(c,t)=\sum_{i=1}^4c_i\otimes\mu_{t,i}
 +\sum_{j,k=1}^3 B(c)_{jk}\otimes(Q_t)_{jk}.
\tag{CAP11}
\]
The first terms are squares by separate square roots.  Factor
$B(c)=UU^T$ within $K_c$ and $Q_t=VV^T$ within $K_t$.  The last
term in (CAP11) is exactly
\[
\sum_{a,b}\left(\sum_jU_{ja}\otimes V_{jb}\right)^2.
\]
This proves the tensor certificate with no mixed denominator.  The
five omitted event squares in (CAP1) are separate positive-source
factors times testing squares and can be restored exactly.  Every
literal projective endpoint is already included in the homogeneous
testing tuple.  Splitting a complex state into its real and imaginary
parts gives the same assertion for complex states.
\end{proof}

\subsubsection*{Every actual finite contact lies strictly below all eight walls}

The fixed event cycle is
\[
(a_0,\ldots,a_7)=(\lambda,D,C,A,\alpha,m,\beta,n).
\]
For indices modulo eight put
\[
\Delta_i=8a_i a_{i+1}a_{i+2}-4d a_i a_{i+2}
 -d^2\{d+2(a_i+a_{i+1}+a_{i+2})\}.
\tag{CAP12}
\]
The cyclic equivalence can be verified without an unspecified symmetry
import.  In the CHSH convention
$S=A_0B_0-A_0B_1+A_1B_0+A_1B_1$, the Bell relabeling
\[
(A'_0,A'_1,B'_0,B'_1)=(B_1,-B_0,-A_0,-A_1)
\tag{CAP13}
\]
preserves $S$ and maps each event in the displayed cycle to the next.
For example the transformed $\lambda=(+,+\mid0,1)$ is the original
$D=(-,+\mid1,1)$.  The complete eight-event permutation and the CHSH
identity are checked exactly.  Relabeling preserves whether each local
pair is noncommuting, so the strict assertion of the first proposition
applies to every cyclic triple.

\begin{theorem}[Automatic contact exclusion]
Let $d>0$, let every event coefficient be nonnegative, and suppose
the original parent (CAP1) has a zero on a nonzero source-field state
at an actual finite noncommuting source pair.  Then
\[
\Delta_i<0\qquad (i=0,\ldots,7).
\tag{CAP14}
\]
\end{theorem}

\begin{proof}
If some $\Delta_i\geq0$, the selected parent for that triple is
positive definite at the specified pair.  The full parent is this
selected parent plus five nonnegative event squares.  It is therefore
positive on every nonzero state at that pair, contradicting the
assumed zero.  This uses only ordinary ordered-field positivity at
the source contact.  It does not evaluate a source contact in an
independent testing field.
\end{proof}

More generally, suppose a proposed cap keeps the same signed coefficient
$d$ and removes nonnegative amounts $e_j$ from event coefficients.
At an actual zero $t_c$ one has the exact identity
\[
F_{\mathrm{cap}}(t_c)
 =-4\sum_j e_j\,\omega_j(t_c)R_j(t_c)^2,
\tag{CAP15}
\]
where $\omega_j$ is the positive angle factor multiplying the event
row in (CAP1).  Thus a globally physical cap can remove a strictly
positive amount only from an event whose row vanishes at that contact.
If one removed positive square has nonzero row, (CAP15) is already a
strictly negative physicality witness.  In the special three-event cap,
(CAP14) rules out even the case in which all removed rows vanish.
The same argument permits smaller retained event coefficients: add
their nonnegative excesses back as well.  It does not apply to a
modification of $d$, whose difference is a signed CHSH term rather
than a nonnegative event square.

\subsubsection*{Exact consequences for source normalization and routing}

For $d>0$ put
\[
\delta_{\rm tot}=\frac{d}{d+\sum_{i=0}^7a_i},
\qquad b_i=a_i/d.
\tag{CAP16}
\]
In a real closed extension, ``bounded'' means bounded in absolute
value by a real number, and ``infinitesimal'' means smaller in
absolute value than every positive real number.  Since all $b_i$
are nonnegative and the sum is finite,
\[
\begin{array}{ll}
\delta_{\rm tot}\text{ has positive real residue}
 &\Longleftrightarrow\text{every }b_i\text{ is bounded},\\
\delta_{\rm tot}\text{ is infinitesimal}
 &\Longleftrightarrow\text{at least one }b_i\text{ is unbounded}.
\end{array}
\tag{CAP17}
\]
These alternatives are disjoint and exhaustive because
$0<\delta_{\rm tot}\leq1$.  They are invariant under a positive
source rescaling and under Bell relabeling.  Normalizing $d=1$ is
legal in the coefficient field, and multiplying each tensor root by
the pure-source factor $\sqrt d$ restores the original scale.

There cannot be three consecutive unbounded $b_i$ at an actual finite
contact.  Indeed division of (CAP12) by its positive triple product
gives
\[
\frac{\Delta_i}{a_i a_{i+1}a_{i+2}}
 =8-\frac{4d}{a_{i+1}}
 -\frac{d^3}{a_i a_{i+1}a_{i+2}}
 -2d^2\left\{
 \frac1{a_i a_{i+1}}+\frac1{a_i a_{i+2}}
 +\frac1{a_{i+1}a_{i+2}}\right\}.
\tag{CAP18}
\]
If the three ratios $a_j/d$ were unbounded, its residue would be
$8$, contradicting (CAP14).  Equivalently, every consecutive triple
contains an event bounded relative to $d$.

If a later chart uses $C=2$, the rescaling factor is $2/C$, and its
new signed weight is $2d/C$; (CAP16) remains unchanged.  The implication
$d/C$ infinitesimal $\Rightarrow\delta_{\rm tot}$ infinitesimal is
valid, but the reverse implication does not follow from normalization
alone.  A different event may be the unbounded one.  An entry map that
claims more must prove its extra source inequalities explicitly.

For a more precise cap screen with $d>0$, (CAP7) yields
\[
b>0\quad\Longrightarrow\quad
\Delta<0\ \Longleftrightarrow\
A<\frac{d^2(d+2C+2D)}{2b}.
\tag{CAP19}
\]
If $b\leq0$, the strict inequality $\Delta<0$ is automatic.  Every
inverse in (CAP19) is within the source field, and the polynomial
condition (CAP14) remains the primary contact condition.



\section{\texorpdfstring{Weak-alpha source entry and the unit-\(n\) case}{Weak-alpha source entry and the unit-n case}}
\label{component:p3a3949869b9e}

\label{p3a3949869b9e:sec:new-weak-alpha-unit-n}

\paragraph{The source domain.}
All coefficients are bounded in the source field, \(d=1\), and the
original canonical source is physical with an actual finite noncommuting
zero contact. The C2 convention is \(\lambda=\beta=0\), and the C1
convention is \(\lambda=n=0\). The first subsection applies to either
family with \(\alpha\) infinitesimal. The all-tests theorem proved below
has the additional C2 hypotheses that \(n\) is a positive unit and \(m\)
is infinitesimal. It includes all independent testing states, signs and
literal projective endpoints. The weak-\(n\) and C1 cases are separate.

\subsubsection*{Source entry, proved before any shape quotient is used}

Physicality forces \(C,A,D\) to be positive units. Here are exact
residue-face tests proving this assertion. In C2 the negative central
coefficient at \(A=0\), and its vanishing positive Laurent terms along
\(U=t^2,V=1/t\), are the displayed forbidden-face proof in the actual
positive-\(d\) source partition. If also \(C=0\) in residue, the source
has only the four alternating events \(n,m,A,D\). At \(u=v=1\) the full
parent at state \((1,-\varepsilon,\varepsilon,1)\) is exactly
\[
 16\varepsilon\{(1+n+m+A+D)\varepsilon-1\}<0
 \tag{WA1}
\]
for a sufficiently small positive real \(\varepsilon\).
Thus \(C\) cannot be weak. The remaining central coefficient identities
are
\[
\begin{aligned}
h|_{\alpha=D=0}
 &=(2A+2C+1)\{8\beta mn-2\beta-4mn-2m-2n-1\},\\
h|_{\alpha=C=n=0}
 &=-(2A+1)(2D+1)(2\beta+2m+1),\\
h|_{\alpha=A=n=0}
 &=-(2C+2D+1)(2\beta+2m+1).
\end{aligned}\tag{WA2}
\]
In the first line \(n\beta=0\), so it is strictly negative. For C2
with \(\alpha=D=0\), the other Laurent terms are
\(q_{20}U/V+q_{21}U\), killed by \(U=t^2,V=t\).
For C1 with \(\alpha=D=0\), they are
\(q_{12}V+q_{21}U+q_{22}UV\), killed by \(U=V=t\).
For C1 with \(\alpha=C=0\), they are \(q_{02}V/U+q_{12}V\),
killed by \(U=1/t,V=t\). For C1 with \(\alpha=A=0\), they are
\(q_{02}V/U+q_{12}V+q_{22}UV\), killed by \(U=t,V=t^2\).
Here \(t\) tends to zero through positive reals. These contradictions
prove all the asserted unit bounds without a source-contact assumption.

Write \(U=u_c^2,V=v_c^2>0\). The actual determinant dictionary now has
\[
 q_{01}=\alpha m a_1,\quad q_{02}=ma_2,\quad
 q_{10}=\alpha a_3,\quad q_{12}=a_4,\quad
 q_{20}=na_5,\quad q_{21}=a_6,\quad q_{22}=\beta a_7,
 \tag{WA3}
\]
where the \(a_i\) are positive bounded units. The two angular equations
and their sum are
\[
\begin{aligned}
 U^2(na_5+a_6V+\beta a_7V^2)&=mV(\alpha a_1+a_2V),\\
 V^2(ma_2+a_4U+\beta a_7U^2)&=\alpha a_3U+na_5U^2,\\
 a_6U+a_4V+2\beta a_7UV
 &=\alpha(ma_1/U+a_3/V).
\end{aligned}\tag{WA4}
\]
Every \(t_i\) in the dictionary is a positive unit because it includes
\(d=1\). At contact all noncentral Laurent terms are nonnegative,
and their sum is bounded. In particular \(U,V,mV/U,nU/V\) are bounded.
Multiplying the last equation by \(V\) shows
\(\alpha/V=O(U+V)\). The first two equations then exclude either
\(U\) or \(V\) having nonzero residue: if \(U\) is a unit the first
forces \(V\) a unit, and conversely the second forces \(U\) a unit;
the last equation then contradicts weak \(\alpha\).
Both are infinitesimal. Consequently
\[
 \alpha/V=o(1),\qquad
 m\asymp (U/V)^2(n+V),\qquad
 \alpha\asymp V(U+V).
 \tag{WA5}
\]
For the last comparison, the denominator
\(a_3+ma_1V/U\) is a positive bounded unit. Actual finite contact
forces \(m,\alpha>0\) in WA4. If \(n\) is a unit and \(m\) weak,
\[
 U/V\asymp\sqrt m=o(1).
 \tag{WA6}
\]
If \(n\) is weak, then, whether \(U/V\) is bounded or unbounded,
WA5 gives \(mV/U,nU/V=o(1)\). These same two conclusions hold
under WA6. They are the common source bounds used in the shape argument.

Orient the actual signs positively, use \(W_c=1\), and put
\[
 p=v_c,\quad q=Z_c,\quad P=p^2,\quad Q=q^2,\qquad
 u_c=rpq,\quad X_c=xpq^2,\quad Y_c=ypq.
 \tag{WA7}
\]
The exact core sign identity SE9 ensures \(W_cZ_c>0\), so this
normalization is legitimate. To prove boundedness before using these
quotients, eliminate the \(X,Y\) block inside the source field. Set
\[
 H_X=n+m+\beta+\tfrac12,\quad K=C+D+\tfrac12,\quad H=K+A .
\]
The block has diagonal positive residue \(H_X,H\) and an infinitesimal
off-diagonal entry. The source Schur entries satisfy
\[
 S_{WW}/U\sim a=AK/H>0,\qquad
 -S_{WZ}/(u_cv_c)\sim b=\tfrac12+AD/H>0.
 \tag{WA8}
\]
Indeed the only additional leading contributions are multiples of
\(mV/U,nU/V,mn\), all infinitesimal on the domains just stated.
The kernel equation gives \(r\sim b/a\), a positive unit. The scaled
\(X,Y\) equations are then a bounded system with positive triangular
residue; its apparent \(m/Q\) term is bounded since
\(m/Q=r^2mV/U\). Thus \(x,y\) are bounded before taking residues.
Those equations and the two remaining kernel equations give
\[
 Ay\sim\tfrac12,\quad
 K(r-y)\sim D+\tfrac12,\quad
 D(1+y-r)\sim r/2,\qquad
 x\sim-\frac{m+\beta+1/2}{H_X}\,r.
 \tag{WA9}
\]
For the last kernel division the possible extra \(\alpha/V\),
\(nU/V\) and \(mV/U\) terms vanish by WA5. In particular \(y,r-y\)
and \(1+y-r\) are positive units, and the critical residue dictionary is
\[
 A_0=\frac1{2y_0},\qquad
 C_0=\frac{1+y_0}{2(r_0-y_0)},\qquad
 D_0=\frac{r_0}{2(1+y_0-r_0)}.
 \tag{WA10}
\]
In the unit-\(n\), weak-\(m\) domain, \(\beta=0\), and hence
\[
 P,Q=o(1),\qquad m\asymp Q^2,\qquad
 \alpha\asymp P^2,\qquad
 x_0=-r_0/(1+2n_0).
 \tag{WA11}
\]
No independent testing coordinate occurred in this entry proof.

\subsubsection*{The exact scalar and every first normal}

For reference the full finite quarter-parent is
\[
\begin{aligned}
F={}&nD_uD_vX^2+m(vW-X+uvY-uZ)^2
       +CD_v(uW-Y)^2+AD_uD_vY^2\\
 &+(D+\tfrac12)(uW+uvX-Y-vZ)^2
       +(\alpha+\tfrac12)D_u(vY-Z)^2\\
 &+(\beta+\tfrac12)D_v(X+uZ)^2-\tfrac12D_uD_vZ^2 .
\end{aligned}\tag{WA12}
\]
Here \(D_u=1+u^2,D_v=1+v^2\). We retain \(\beta\) symbolically to make
the exact common graph explicit, but the current case has \(\beta=0\).
In the testing chart \(W=1,Z=z\), use
\[
 u_g=rvz,\qquad X_g=xvz^2,\qquad Y_g=yvz,\qquad
 V_t=v^2,\quad T=z^2.
 \tag{WA13}
\]
In the rest of this subsection write \(V=V_t\) only. The full graph
restriction is
\[
\begin{aligned}
\Phi={}&mV+\alpha T+hVT+\mathcal C V^2T+EVT^2
   +\mathcal F V^2T^2+GV^3T^2+JV^2T^3+\mathcal K V^3T^3,\\
\mathcal C={}&Ay^2+C(r-y)^2+(\alpha+\tfrac12)y^2+2mry,\\
E={}&\alpha r^2+(m+\beta+\tfrac12)(r+x)^2+nx^2,\\
G={}&r^2y^2(A+\alpha+m+\tfrac12),\qquad
J=nr^2x^2,\qquad \mathcal K=(D+n+\tfrac12)r^2x^2,\\
\mathcal F={}&Ar^2y^2+(2D+1)rx(r-y-1)-(2\alpha+1)r^2y\\
 &-2mry(r+x)+\beta(r+x)^2+nx^2+rx+x^2/2 .
\end{aligned}\tag{WA14}
\]
The coefficient \(h\) is defined by this entire expansion; it is bounded.
Every displayed coefficient is bounded. In the current unit-\(n\)
domain \(\mathcal C,E\) are positive units, since
\[
 \mathcal C_0=r_0(1+y_0)/2,\qquad
 E_0=n_0r_0^2/(1+2n_0).
 \tag{WA15}
\]
The coefficient \(\mathcal F\) need not be positive.

The actual contact gives \(\Phi(P,Q)=\Phi_V(P,Q)=\Phi_T(P,Q)=0\).
Solving these three linear equations in \(m,\alpha,h\) is exact:
\[
\begin{aligned}
m={}&EQ^2+\mathcal F PQ^2+GP^2Q^2+2JPQ^3+2\mathcal K P^2Q^3,\\
\alpha={}&\mathcal C P^2+\mathcal F P^2Q+2GP^3Q
                         +JP^2Q^2+2\mathcal K P^3Q^2,\\
h={}&-2\mathcal CP-2EQ-3\mathcal FPQ-4GP^2Q
                         -4JPQ^2-5\mathcal K P^2Q^2 .
\end{aligned}\tag{WA16}
\]
Put
\[
 R=z(V-P),\quad S=v(T-Q),\quad
 \Psi=QR^2+PS^2+vzRS .
\]
With the exact balanced monomials
\[
\begin{aligned}
B_{32}&=2PQR^2+(vzR+PS)^2,\\
B_{23}&=2PQS^2+(QR+vzS)^2,\\
B_{33}&=(2PQ^2+VT^2)R^2+(2P^2Q+PVT)S^2
                              +(3PQ+PT)vzRS,
\end{aligned}\tag{WA17}
\]
the whole scalar, not a leading truncation, is
\[
 \Phi=\mathcal CR^2+ES^2+\mathcal F\Psi
                          +GB_{32}+JB_{23}+\mathcal K B_{33}.
 \tag{WA18}
\]
For \(v,z\) infinitesimal, all entries after the first two positive
unit pivots are infinitesimal, by \(P,Q=o(1)\). This is a strict
two-row Gram even when \(\mathcal F\) is signed.

Let \(L_u,L_X,L_Y\) be half the full first derivatives of WA12 at
WA13. The complete polynomials have the forms
\[
\begin{aligned}
 L_u&=vz A_u(V,T),\qquad L_Y=vz A_y(V,T),\\
 A_u&=a_{00}+a_{10}V+a_{01}T+a_{11}VT
                      +a_{12}VT^2+a_{21}V^2T+a_{22}V^2T^2,\\
 A_y&=e_{00}+e_{10}V+e_{11}VT+e_{21}V^2T,\\
 L_X&=v[-m+TB_X(V,T)],\qquad
 B_X=c_0+c_1V+c_2VT+c_3V^2T .
\end{aligned}\tag{WA19}
\]
All coefficients are bounded, and the expansion fixes every one
unambiguously. In particular
\[
 c_2=nr^2x,\qquad c_3=(D+n+\tfrac12)r^2x .
\]
Their actual source jets are \(A_u(P,Q)=A_y(P,Q)=0\) and
\(B_X(P,Q)=m/Q\). Finite monic difference division gives
\[
\begin{aligned}
 L_i&=v A_{i,V}(V,T)R+z B_i(T)S,\qquad i=u,Y,\\
 L_X&=vz B_{X,V}(V,T)R+
                \{T(Pc_2+P^2c_3)+m/Q\}S .
\end{aligned}\tag{WA20}
\]
Here \(A_{i,V}=[A_i(V,T)-A_i(P,T)]/(V-P)\), and \(B_{X,V}\)
is defined identically. The other divided polynomials are
\[
\begin{aligned}
B_u(T)&=a_{01}+Pa_{11}+Pa_{12}(T+Q)+P^2a_{21}
                                      +P^2a_{22}(T+Q),\\
B_Y(T)&=Pe_{11}+P^2e_{21}.
\end{aligned}
\]
These are literal polynomial identities with bounded coefficients,
including \(v=0\) or \(z=0\). Since \(m/Q=O(Q)=o(1)\),
every full first normal is an infinitesimal linear combination of \(R,S\).

Set \(N_u=u-rvz,N_X=X-xvz^2,N_Y=Y-yvz\). In the chart under
consideration these are infinitesimal. The full higher-normal
remainder is a polynomial in the three normals with no terms of
normal degree below two; all terms of degree at least three give
infinitesimal Gram entries. Its quadratic residue is
\[
 K_0(N_u-N_Y)^2+A_0N_Y^2+(n_0+\tfrac12)N_X^2\succ0.
 \tag{WA21}
\]
WA18 and WA20 therefore give a strict bounded Gram for the entire
parent on the three hard and two soft rows. The strict-Gram lemma
proves its tensor-SOS assertion, with no mixed inverse.

\subsubsection*{The complementary local states}

Normalize a maximal testing state coordinate. At \(u,v=o(1)\)
the state residue has positive pivots on \(X,Y\); if one is a unit,
the full parent has positive residue immediately. Otherwise one
of \(W,Z\) is a unit. The preceding proof covers \(W=1,Z=o(1)\).
Every remaining state can be normalized to \(Z=1,W=\omega\)
bounded. Retain the entire positive \(m,\alpha\) squares.
The remainder lies exactly in the square of \((u,v,X,Y)\).
Its quadratic real residue is
\[
\begin{aligned}
 &(n_0+\tfrac12)\left(X+\frac{u}{1+2n_0}\right)^2
     +\frac{n_0}{1+2n_0}u^2\\
 &+H_0\left(Y-\frac{K_0\omega u-D_0v}{H_0}\right)^2
     +a_0(\omega u-r_0v)^2,\qquad a_0=A_0K_0/H_0>0 .
\end{aligned}\tag{WA22}
\]
The critical identities WA10 prove the equality. This quadratic
form is positive definite for every real \(\omega\), including zero.
All higher normal terms are infinitesimal corrections on these
four rows, while \(\omega-\bar\omega\) is retained in coefficients.
Thus every state at the double-infinitesimal angle corner is covered.

\subsubsection*{All other testing angles, including the literal endpoints}

Again retain the entire \(m,\alpha\) squares and call the remainder
\(F_0\). Its source residue is the critical \(C,A,D,n\) source.
For bounded \(u,v\), \(F_0\) is exactly a bounded quadratic form
in the four rows \((uW,X,Y,Z)\). Its real \(Z=0\) block is positive:
the \(nX^2\) and \(AY^2\) squares and \(C(uW-Y)^2\) suffice.
The determinant is
\[
 \det \operatorname{Mat}_{(uW,X,Y,Z)}(\bar F_0)
 =D_u^2D_v^2\{q_{20}u^2+q_{12}v^4+q_{21}u^2v^2\}/16>0
 \tag{WA23}
\]
whenever at least one of the two real angle residues is nonzero.
Here \(q_{20},q_{12},q_{21}>0\) are the actual residue determinant
coefficients, and \(h=0\) follows from WA10. This proves strictness
of the full four-row Gram. Literal \(u=0\) is allowed.

For Bob reciprocal \(b=1/v\) infinitesimal and Alice bounded,
use the same rows \((uW,X,Y,Z)\) in \(b^2F_0(u,1/b)\).
At \(b=0\) its whole real expression is
\[
 C_0(uW-Y)^2+(A_0+\tfrac12)D_uY^2
             +D_0(uX-Z)^2+(n_0+\tfrac12)D_uX^2\succ0 .
 \tag{WA24}
\]
It is strict at every real \(u\), including zero.
For Alice reciprocal \(s=1/u\) infinitesimal and Bob bounded,
the original state rows suffice. The expression at \(s=0\) is
\[
 n_0D_vX^2+C_0D_vW^2+A_0D_vY^2
          +(D_0+\tfrac12)(W+vX)^2+\tfrac12(vY-Z)^2\succ0 .
 \tag{WA25}
\]
This includes \(v=0\), using the positive \(n_0\) pivot.

For both reciprocal angles \(s,b=o(1)\), use the five rows
\((W,X,Y,sZ,bZ)\). The whole reciprocal polynomial belongs
to their square ideal, and its real coefficient Gram is
\[
 C_0W^2+A_0Y^2+D_0(X-sZ)^2
          +(n_0+\tfrac12)X^2+\tfrac12(Y-bZ)^2\succ0 .
 \tag{WA26}
\]
Every remaining coefficient is infinitesimal. These exact finite
square-ideal allocations include \(s=0\), \(b=0\) and arbitrary
states; no projective endpoint is obtained by a limiting argument.

The four angle charts and the two local state charts are exhaustive.
Testing sign gauges, positive pure-field homogeneous normalization
factors and real/imaginary state splitting restore the original parent:
\[
 F_c(u,v;\psi)\ \in\ \Sigma(K_c\otimes_{\mathbb R}K_t)^2
 \quad\text{for every independent projective test and state.}
 \tag{WA27}
\]
This closes exactly the C2 unit-\(n\), weak-\(m,\alpha\) case.
The common source-entry subsection does not by itself close the
weak-\(n\) or C1 weak-\(\alpha\) branches.



\section{\texorpdfstring{Weak-alpha weak-\(n\) and C1: the complete parent}{Weak-alpha weak-n and C1: the complete parent}}
\label{component:p0db05df21ad1}

\label{p0db05df21ad1:sec:weak-alpha-weak-n-c1}

\paragraph{The source domain.}
The source is an actual physical canonical finite noncommuting contact,
all coefficients are bounded in the source field, and \(d=1\).
Either it is C2 with \(\lambda=\beta=0\) and \(n,\alpha\) infinitesimal,
or it is C1 with \(\lambda=n=0\), \(\alpha\) infinitesimal and arbitrary
bounded \(\beta\geq0\). All strict source gaps retain their original
meaning. We prove the tensor-SOS statement for the entire parent at
every independent projective testing angle and every state.
The common source-entry proof WA1--WA10 is used on precisely its
weak-\(n\) domain. In particular \(C,A,D,r,y,r-y,1+y-r\) are positive
units, \(x\) is bounded, and \(x_0=-r_0\).
Boundedness, infinitesimality and residue refer
to the separated bounded tensor ring and its real residue map.
The strict-Gram lemma is applied only after a finite explicit row
allocation; no mixed tensor is inverted.

Use the full quarter-parent WA12 and its exact source coordinates
\[
 p=v_c,\quad q=Z_c>0,\quad P=p^2,\quad Q=q^2,\qquad
 W_c=1,\quad u_c=rpq,\quad X_c=xpq^2,\quad Y_c=ypq.
 \tag{WN1}
\]
Thus \(P,PQ\) are positive infinitesimals. Put
\[
 H_0=m+\beta+\tfrac12,\quad H_X=H_0+n,\quad
 K=C+D+\tfrac12,\quad H_Y=K+A.
 \tag{WN2}
\]
These are positive units. The exact scalar coefficients
\(\mathcal C,E,\mathcal F,G,J,\mathcal K\), the contact balances, and
the complete first normal polynomials are those displayed in
WA14, WA16 and WA19; those identities retain \(\beta\) and \(n\).
Their relevant residues are
\[
 \mathcal C_0=\frac{r_0(1+y_0+4m_0y_0)}2>0,\qquad
 \mathcal F_0=\frac{r_0^2(1+y_0)}2>0,
 \qquad G,J,\mathcal K\geq0 .
 \tag{WN3}
\]
Both \(\mathcal C,\mathcal F\) are consequently positive units.
The exact source \(X\)-jet gives, with the WA19 coefficients,
\[
\begin{gathered}
 c_0=H_0(r+x)+nx,\quad c_2=nr^2x,\quad
 c_3=(D+n+\tfrac12)r^2x,\\
 c_0=m/Q-Pc_1-PQc_2-P^2Qc_3,\\
 E=\alpha r^2+c_0^2/H_X+nH_0r^2/H_X .
\end{gathered}\tag{WN4}
\]
The last equality is an exact completion inside the source field.

\subsubsection*{Two finite positive circuit identities}

For \(V=v^2,T=z^2\) put
\[
 R=z(V-P),\quad S=v(T-Q),\quad
 \Psi=V^2T^2+PQ^2V+P^2QT-3PQVT .
\]
Then \(\Psi=QR^2+PS^2+vzRS\). It is an explicit six-square
cubic circuit: if \(L_1^3=V^2T^2,L_2^3=PQ^2V,L_3^3=P^2QT\),
then \(L_1L_2L_3=PQVT\), and
\[
 \Psi=\sum_{i=1}^3 L_i
 \bigl\{(L_1-L_2)^2-(L_1-L_2)(L_1-L_3)+(L_1-L_3)^2\bigr\}.
 \tag{WN5}
\]
The real \(2\)-by-\(2\) Gram in each brace is strictly positive.
All radicals in this formula are elementary separated monomials:
source radicals and testing radicals are taken in their own real
closed fields and then tensored.

The balanced degree-\((3,3)\) term has the polynomial identity
\[
 a^{10}+2b^{10}+2c^{10}-5a^2b^4c^4
 =2(b^5-c^5)^2
 +(a^2-bc)^2(a^6+2a^4bc+3a^2b^2c^2+4b^3c^3).
 \tag{WN6}
\]
For nonnegative elementary monomials \(a,b,c\), the bracket is the
sum of squares of the four rows
\[
 a^3,\qquad \sqrt2a^2\sqrt{bc},\qquad
 \sqrt3abc,\qquad 2(bc)^{3/2}.
\]
This is a five-square separated certificate, including vanishing
variables. Applied with
\[
 a^{10}=V^3T^3,\quad b^{10}=P^2Q^3V,\quad c^{10}=P^3Q^2T,
\]
it proves positivity of the entire \(B_{33}\) in WA17.

A further exact identity supplies any fixed positive real pivot
on the row \(PS\):
\[
 2P\Psi+VT R^2-(PS)^2=(PS+vzR)^2+2PQR^2.
 \tag{WN7}
\]
Indeed, for every fixed \(b>0\),
\[
\begin{split}
 \mathcal C R^2+\mathcal F\Psi={}&
 (\mathcal C-bVT+2bPQ)R^2+(\mathcal F-2bP)\Psi\\
 &+b(PS)^2+b(PS+vzR)^2.
\end{split}\tag{WN8}
\]
On the charts where this is used, \(VT,P,PQ\) are infinitesimal,
so the original two unit residues are unchanged.

\subsubsection*{Bounded source ratio \(Q\)}

From WA5, WN4 and bounded \(Q\),
\[
 m\asymp Q^2(n+P),\quad \alpha\asymp P^2,\quad
 r+x=O(n+P),\quad E\asymp n+P^2.
 \tag{WN9}
\]
In particular \(m,n,E,\alpha\) are infinitesimal. None of these
estimates contains an independent testing coordinate.

First take \(W=1,Z=z\) bounded and \(u,v\) infinitesimal. Use the
entire graph \(u=rvz,X=xvz^2,Y=yvz\). The exact scalar is
\[
 \Phi=\mathcal CR^2+ES^2+\mathcal F\Psi
          +GB_{32}+JB_{23}+\mathcal K B_{33},
 \tag{WN10}
\]
where the last three terms retain their positive full certificates.
The strict soft rows are \(R\), the six cubic rows in WN5,
\(\sqrt E\,S\), and \(PS\); WN8 supplies any needed real pivot
on the last row.

The exact monic normal divisions WA20 apply unchanged.
Their \(R\)-coefficients contain \(v\) or \(vz\), hence are
infinitesimal. Their \(S\)-coefficients are \(O(E+P)\):
in particular
\[
 a_{01}=\alpha r+H_0(r+x)=O(E+P),\qquad
 a_{12}=nr x^2,\qquad m/Q=O(E+P).
 \tag{WN11}
\]
Every \(ES\) contribution has infinitesimal coefficient \(\sqrt E\)
on \(\sqrt E S\). Coefficients on \(PS\) may be bounded units.
There are only finitely many of them, so choose \(b\) in WN8
larger than their finite real Schur cost. This uses a real constant,
not an inverse of a mixed tensor.

The three hard normals
\(N_u=u-rvz,N_X=X-xvz^2,N_Y=Y-yvz\) are infinitesimal.
Their complete quadratic residue is
\[
 K_0(N_u-N_Y)^2+A_0N_Y^2+
          (\beta_0+\tfrac12)(N_X+\bar zN_u)^2\succ0.
 \tag{WN12}
\]
All higher-normal terms have degree at least three and bounded
coefficients, hence allocate as infinitesimal Gram entries on these
hard rows. WN10--WN12 prove the full parent in this chart.

For the complementary local states take \(Z=1,W=w\) infinitesimal,
set \(a=u/r\), and use
\[
 v=aw,\quad X=xa,\quad Y=yaw,\qquad U=a^2,\quad T=w^2.
 \tag{WN13}
\]
The entire scalar is
\[
\begin{split}
 \Phi'={}&\alpha+mUT^2+EU+hUT+\mathcal C U^2T^2
       +\mathcal F U^2T+GU^3T^2+JU^2+\mathcal K U^3T .
\end{split}\tag{WN14}
\]
Here \(a=(1/r)\otimes u\) is elementary; its roots remain separated.
The strict rows are \(P,a^2w,\sqrt E\,a\).
The term \(hUT\) is harmless: the contact balances give
\(h=O(P+E)\), and writing \(h=Ph_P+Eh_E\) with bounded source
coefficients gives the infinitesimal Gram terms
\(h_Pw\,P(a^2w)+h_Ew^2(\sqrt E a)^2\).
All other terms of WN14 are positive.

For completeness, half the entire \(v,Y,X\) derivatives have
the polynomial forms
\[
\begin{aligned}
 L_v'&=aw(d_0+mT+d_1U+d_2UT+d_3U^2+d_4U^2T),\\
 L_Y'&=aw(e_0+e_1U+e_2UT+e_3U^2T),\\
 L_X'&=a(c_0-mT+c_2U+c_1UT+c_3U^2T).
\end{aligned}\tag{WN15}
\]
Every coefficient is bounded. Source jets at
\(U=PQ,T=1/Q\) give
\[
\begin{split}
d_0&=-m/Q-d_1PQ-d_2P-d_3P^2Q^2-d_4P^2Q=O(n+P),\\
e_0&=-e_1PQ-e_2P-e_3P^2Q=O(P).
\end{split}
\]
Also \(c_0,m=O(n+P)\) and \(c_2=O(n)\).
Splitting \(n+P=O(E+P)\) in these displayed complete polynomials
allocates all first normals with infinitesimal coefficients to
\(P,a^2w,\sqrt E a\). For example \(awP\), \(awE\), and
\(a^3w\) have coefficients \(aw\), \(w\sqrt E\), and \(a\)
on those three rows; the other monomials have still more small
testing factors. The complete hard residue is
\[
 (\beta_0+\tfrac12)N_X^2+
 D_0N_v^2+2D_0N_vN_Y+H_{Y,0}N_Y^2\succ0,
 \tag{WN16}
\]
where \(N_v=v-aw,N_Y=Y-yaw,N_X=X-xa\).
The higher-normal argument is again exact polynomial allocation.
This proves all local states when \(Q\) is bounded.

\subsubsection*{Unbounded source ratio: the \(W\)-chart}

Suppose \(Q\) is unbounded. The source estimates and positive
contact balances give
\[
 PQ=o(1),\quad m\asymp Q^2(n+P),\quad nQ^2,PQ^2=O(1),
 \quad \alpha\asymp P^2Q .
 \tag{WN17}
\]
WN4 now gives
\[
 c_0\sim m/Q,\quad r+x\sim m/(H_XQ),\quad
 E\asymp n+(m/Q)^2,\quad h/m=o(1).
 \tag{WN18}
\]
To justify the first comparison before its use, divide the three
errors \(Pc_1,PQc_2,P^2Qc_3\) by \(m/Q\).
Their absolute values are \(O(1/Q),O(P),O(P)\), respectively;
\(n/(m/Q)=O(1/Q)\). The positive term
\(\alpha/(m/Q)^2=O(1/Q)\) follows from
\(m\geq cPQ^2\). The last comparison follows term by term from
WA16 using these same inequalities.

Take \(W=1,Z=z\) bounded, \(u,v\) infinitesimal. If \(m\) is a
unit, retain the entire \(n,\alpha\) squares. The remainder lies
in the square of \((u,v,X,Y)\). Its real quadratic form is
\[
\begin{split}
 &H_0\{X+zu-(m/H_0)v\}^2+
       \frac{m(\beta+1/2)}{H_0}v^2\\
 &\qquad+H_Y\{Y-(Ku-Dzv)/H_Y\}^2
       +a_*(u-rzv)^2,\qquad a_*=AK/H_Y>0,
\end{split}\tag{WN19}
\]
with coefficients interpreted in residue. WA10 verifies this
identity. It is strict for every bounded \(z\), so the complete
higher terms give a strict Gram.

If \(m\) is weak, use the graph WA13. Its entire scalar has strict
rows \(\sqrt m\,v,v^2z,\sqrt\alpha\,z\), since
\(mV+hVT=mV[1+(h/m)T]\), while \(\mathcal C,\mathcal F,E,G,J,\mathcal K\)
are nonnegative. The complete source jet equations yield
\[
 a_{00}=O(m),\quad e_{00}=O(PQ),\quad
 a_{01},c_0=O(m/Q).
 \tag{WN20}
\]
Thus \(a_{00}/\sqrt m\), \(e_{00}/\sqrt m\),
\(a_{01}/\sqrt m\), \(c_0/\sqrt m\) are infinitesimal.
For \(e_{00}\), use \(PQ/\sqrt m=O(\sqrt P)\).
Every other monomial of the complete normals WA19 has a factor
\(v^3z\) or \(v^3z^2\), so it has an infinitesimal coefficient
on \(v^2z\). The hard residue is WN12; all higher terms are
allocated exactly as there. This proves the whole \(W\)-chart.

\subsubsection*{Unbounded source ratio: the complementary scalar}

Use WN13 with \(a,w\) infinitesimal. Introduce only source quantities
\[
\begin{gathered}
 \epsilon=Q^{-1}=o(1),\quad \Lambda=PQ=o(1),\quad
 \tau=\Lambda/\epsilon=PQ^2=O(1),\\
 \nu=n/\epsilon^2=O(1),\qquad N=E/\epsilon^2=O(1).
\end{gathered}\tag{WN21}
\]
The quantities \(\epsilon,\Lambda,\tau,N\) are positive; \(\nu\geq0\).
Let
\[
 A_1=a(T-\epsilon),\quad K_1=UT-\Lambda\epsilon,\quad
 \mathcal M=U^2T+\tau UT^2+\Lambda^2\epsilon-3\Lambda UT .
 \tag{WN22}
\]
The whole scalar WN14, using exactly the three contact balances, is
\[
\begin{split}
 \Phi'={}&\mathcal C K_1^2+NA_1^2+\mathcal F\mathcal M\\
 &+G\{2\Lambda K_1^2+UT^2(U-\Lambda)^2\}\\
 &+J\{(U-\Lambda)^2+2\Lambda A_1^2/\epsilon^2\}
       +\mathcal K\mathcal B_{33},\\
 \mathcal B_{33}={}&U^3T+2(\Lambda^2/\epsilon)UT^2
                      +2\Lambda^3\epsilon-5\Lambda^2UT .
\end{split}\tag{WN23}
\]
The apparent \(1/\epsilon^2\) in the positive \(J\) square is
a source-field quantity, and \(J/\epsilon^2=\nu r^2x^2\) is bounded.
All positive terms are retained. WN6 applies to \(\mathcal B_{33}\)
with \(a^{10}=U^3T,b^{10}=(\Lambda^2/\epsilon)UT^2,
c^{10}=\Lambda^3\epsilon\).
Here \(\Lambda^2/\epsilon=\tau\Lambda\) is bounded.

Write \(a=t^3,w=s^3,\epsilon=h^6,\tau=c^6\), taking positive
elementary roots. Then \(t,s,h\) are infinitesimal and \(c\)
is bounded; it is allowed to be a unit. Put
\[
\begin{gathered}
 L_1=t^4s^2,\quad L_2=c^2t^2s^4,\quad L_3=h^6c^4,\\
 f_1=t^2s,\quad f_2=cts^2,\quad f_3=h^3c^2,\\
 R_{2i-1}=f_i(L_1-L_2),\qquad R_{2i}=f_i(L_1-L_3)
 \quad (i=1,2,3).
\end{gathered}\tag{WN24}
\]
Then \(\mathcal M=\sum_{i=1}^3(R_{2i-1}^2-R_{2i-1}R_{2i}+R_{2i}^2)\),
a strict six-row constant Gram. The following are exact identities:
\[
\begin{aligned}
 aw(U-\Lambda)&=t^3R_1+c^2ts^2R_2,\\
 a^2w(T-\epsilon)&=-h^6R_1+s^6R_2-h^3s^3R_5,\\
 \tau aw(T-\epsilon)&=c^2ts^2(R_2-R_1)-c^3s^3R_3,\\
 \tau\sqrt\epsilon\,A_1&=t^3R_6-t(t^2+c^2s^2)R_5 .
\end{aligned}\tag{WN25}
\]
Every coefficient on the right is infinitesimal, even when \(c\)
is a unit. These are finite linear conductors, not limiting
approximations. The further quadratic conductor is
\[
 \boxed{\ \Lambda(U-\Lambda)K_1
       =\Lambda\mathcal M-(\tau\sqrt\epsilon\,A_1)^2.\ }
 \tag{WN26}
\]
Its right side is an infinitesimal Gram on the six cubic rows.
It remains so after multiplication by any bounded signed coefficient.

If \(\tau\) is a unit, then \(K_1\) itself has an infinitesimal
linear representation on those rows:
\[
 K_1=\tau^{-1/3}
       \{\sqrt{L_1}(R_2-R_1)+\sqrt{L_3}R_6\}.
 \tag{WN27}
\]
This follows by expansion as
\(\tau^{-1/3}(L_1L_2-L_3^2)=UT-\Lambda\epsilon\).
Consequently any fixed real pivot on \(K_1\) can be supplied at
infinitesimal cost to the strict \(\mathcal M\) Gram in this case.

\subsubsection*{Every complementary first normal}

Let \(D_U,E_U\) denote the monic \(U\)-divided differences of
the bracket polynomials in the first two lines of WN15.
The source jets at \(U=\Lambda,T=\epsilon\) give the exact
identities
\[
\begin{aligned}
 L_v'={}&aw\{D_U(U,T)(U-\Lambda)
                 +(m+\Lambda d_2+\Lambda^2d_4)(T-\epsilon)\},\\
 L_Y'={}&aw\{E_U(U,T)(U-\Lambda)
                 +(\Lambda e_2+\Lambda^2e_3)(T-\epsilon)\},\\
 L_X'={}&-mA_1+
       a(c_2+c_3\Lambda\epsilon)(U-\Lambda)
                  +a(c_1+c_3U)K_1 .
\end{aligned}\tag{WN28}
\]
There are no omitted monomials. Define
\[
 d_*=(c_2+c_3\Lambda\epsilon)/\epsilon
       =\nu\epsilon r^2x+\Lambda c_3 .
\]
Since \(\epsilon(U-\Lambda)=K_1-U(T-\epsilon)\), the last line is
\[
 L_X'=-mA_1+P_X,\qquad
 P_X=a(c_1+c_3U+d_*)K_1-d_*UA_1 .
 \tag{WN29}
\]
The exact source equations also give
\[
\begin{aligned}
 m&=N+\tau(\mathcal F+G\Lambda\epsilon+2\mathcal K\Lambda)
                            +2\nu r^2x^2\Lambda,\\
 c_0/\epsilon&=m-\Lambda B,\qquad
 B=c_1+\Lambda c_3+\nu\epsilon r^2x,\\
 N&=\alpha r^2/\epsilon^2+
             (m-\Lambda B)^2/H_X+\nu H_0r^2/H_X .
\end{aligned}\tag{WN30}
\]
In particular
\[
 N,\nu,\tau=O(m),\quad \Lambda=O(m\epsilon),\quad
 N=m^2/H_X+\nu H_0r^2/H_X+O(m\Lambda).
 \tag{WN31}
\]
For the error term use \(\alpha/\epsilon^2=O(\tau\Lambda)\),
which is the exact positive alpha balance WN23/WA16.
All these comparisons are inside the source field.
We split by the bounded source ratio \(\nu/m\); \(m>0\).

\paragraph{Substantial \(\nu/m\).}
If \(\nu/m\) has positive residue, keep \(S_0=\sqrt m A_1\).
The \(N_X,S_0\) Gram, including the entire \(-mA_1\) first
normal, has residue from
\[
 \begin{pmatrix}H_X&-\sqrt m\\-\sqrt m&N/m\end{pmatrix}.
 \tag{WN32}
\]
Its determinant is
\((\nu/m)H_0r^2+O(\Lambda)\), a positive unit.
The \(v,Y\) first normals in WN28 use the first conductor WN25;
their remaining coefficients on \(S_0\) are
\(O(\sqrt m\,w)\), hence infinitesimal.
In WN29 the \(K_1\)-coefficient contains \(a\), and
\(d_*U/\sqrt m=O(\sqrt m\,\epsilon U)\) is infinitesimal.
Thus all other first normals have infinitesimal coefficients
on \(K_1\), the six cubic rows, and \(S_0\).
The \(v,Y\) hard block in WN16 is strict. All higher terms have
at least three infinitesimal hard normals and bounded coefficients.
This proves the full parent on this source subdomain.

\paragraph{Thin \(\nu/m\), including literal \(n=0\).}
Assume \(\nu/m=o(1)\). Equations WN30--WN31 imply
\(\tau/m\) is a positive unit: after division by \(m\),
\[
 1=N/m+(\tau/m)\mathcal F+o(1),\qquad
 N/m=m/H_X+o(1),
\]
and \(1-m/H_X=(n+\beta+1/2)/H_X\) has positive residue.
Write
\[
 N-m^2/H_X=\nu\xi+\eta,\qquad
 \xi=H_0r^2/H_X>0,\qquad \eta=O(m\Lambda).
\]
Complete only the source coefficient \(H_X\):
\[
\begin{split}
 H_XN_X^2-2mN_XA_1+NA_1^2
 &=H_X\widetilde N_X^2+\nu\xi A_1^2+\eta A_1^2,\\
 \widetilde N_X&=N_X-(m/H_X)A_1.
\end{split}\tag{WN33}
\]
Retain the strict row \(\sqrt\nu A_1\) when \(\nu>0\);
if \(\nu=0\), all terms allocated to it vanish literally.
The signed error is an infinitesimal cubic Gram, since
\[
 m\Lambda A_1^2=(m/\tau)(\tau\sqrt\epsilon\,A_1)^2.
 \tag{WN34}
\]
The first two normals in WN28 use WN25 and bounded \(m/\tau\).
For \(P_X\), its \(K_1\)-coefficient contains \(a\), while its
two other contributions have infinitesimal allocations
\[
 \nu\epsilon U A_1=(\sqrt\nu\,\epsilon U)(\sqrt\nu A_1),
 \qquad \Lambda U A_1=(\sqrt\epsilon\,U)
                         (\tau\sqrt\epsilon\,A_1).
 \tag{WN35}
\]
Therefore \(2P_X\widetilde N_X\) has infinitesimal Gram entries.

The shift in WN33 also creates the entire scalar term
\(2(m/H_X)A_1P_X\), which must not be dropped.
Its main factor is controlled by the exact identity
\[
 mU(T-\epsilon)K_1
   =mK_1^2-(m/\tau)\Lambda(U-\Lambda)K_1 .
 \tag{WN36}
\]
The second term is infinitesimal by WN26. If \(m\) is
infinitesimal, the first is an infinitesimal \(K_1\)-Gram.
If \(m\) is a unit, then \(\tau\) is a unit and WN27
makes it an infinitesimal cubic Gram. In that case a positive
real \(K_1\)-pivot may be reinserted by WN27 at infinitesimal
cubic cost. This allocation remains valid with each bounded
factor \(c_1+c_3U+d_*\).
The remaining scalar factors are multiples of
\(m\nu\epsilon U A_1^2\) and \(m\Lambda U A_1^2\);
they have infinitesimal Gram entries on \(\sqrt\nu A_1\)
and the fourth conductor in WN25, respectively.
This accounts for the whole scalar effect of the completion.

\subsubsection*{The entire higher-normal restoration in the thin case}

The full parent is quadratic in \(X\). Its coefficient beyond
the source constant \(H_X\) is exactly the nonnegative polynomial
\[
 D_X=(n+\beta+\tfrac12)v^2+nu^2+
                         (n+D+\tfrac12)u^2v^2 .
 \tag{WN37}
\]
Retain the entire square \(D_XN_X^2\), evaluated at the actual
independent \(u,v\), as separated squares. Each monomial coefficient
is a nonnegative source scalar times a testing square. It remains
a positive whole square after WN33.
The remaining polynomial has no quadratic \(N_X\) term except
\(H_XN_X^2\). At the complementary graph its half mixed
derivatives have the exact forms
\[
 [N_XN_v]=w[-m+a^2B_v],\qquad
 [N_XN_Y]=a^2wB_Y,
 \tag{WN38}
\]
where \(B_v,B_Y\) are bounded polynomials fixed by the full
parent expansion. After the shift, the new first-normal terms
are bounded multiples of \(\tau aw(T-\epsilon)\), using
\(m/\tau=O(1)\); WN25 gives infinitesimal coefficients.
Every other remaining monomial linear in \(N_X\) has degree at
least two in the other hard normals. Replacing \(N_X\) by
\(\widetilde N_X+(m/H_X)A_1\) gives an infinitesimal coefficient
on a product of at least two hard normals, since \(A_1=o(1)\).
Terms involving \(\widetilde N_X\) have total hard-normal
degree at least three and are handled in the same way.
Terms independent of \(N_X\) are unchanged. The full real hard block is
\[
 H_{X,0}\widetilde N_X^2+
       D_0N_v^2+2D_0N_vN_Y+H_{Y,0}N_Y^2\succ0.
 \tag{WN39}
\]
Equations WN23--WN39 therefore give a finite strict Gram for
the whole parent, including all first normals, every higher term,
all \(\beta\)-openings and literal \(n=0\).
This completes the double-infinitesimal testing corner for both
source ratios and all local states.

\subsubsection*{Every other angle, all states and literal endpoints}

Retain the entire \(m,n,\alpha,\beta\) squares and call the
remainder \(F_0\). Only its residue is the critical \(C,A,D\)
source in WA10; no global physicality of the subtracted finite
source is asserted. The following are complete polynomial row
allocations, so their strict residue Grams apply to the original
nearby coefficients.

For bounded \(u,v\) with \(v\) a unit use \((uW,X,Y,Z)\).
The \(Z=0\) principal block is positive definite and
\[
 \det\operatorname{Mat}_{(uW,X,Y,Z)}\bar F_0
 =D_u^2D_v^2(q_{12}v^4+q_{21}u^2v^2)/16>0.
 \tag{WN40}
\]
Both \(q_{12},q_{21}\) are positive units. If \(u\) is a unit
and \(v\) infinitesimal, use instead
\((uW-Y,Y,X+uZ,vZ)\). The entire parent is polynomial in
these rows with bounded coefficients; the only algebraic
coefficient divisions are by the pure testing unit \(u\).
Its determinant at \(v=0\) is
\(D_u^2q_{21}u^2/16>0\), and the first three-row principal block
is strict. This covers all bounded angle pairs outside the
double-infinitesimal corner.

For Bob reciprocal \(b=1/v=o(1)\), Alice bounded, the rows
\((uW,X,Y,Z)\) give the real form
\[
 C_0(uW-Y)^2+(A_0+\tfrac12)D_uY^2+
                 D_0(uX-Z)^2+\tfrac12D_uX^2\succ0.
 \tag{WN41}
\]
For Alice reciprocal \(s=1/u=o(1)\) and Bob a bounded unit,
the original state rows give
\[
 C_0D_vW^2+A_0D_vY^2+
       (D_0+\tfrac12)(W+vX)^2+\tfrac12(vY-Z)^2\succ0.
\]
If \(s,v\) are both infinitesimal, use
\((W,Y,vX,\Gamma)\), where \(\Gamma=sX+Z\).
The entire reciprocal parent is polynomial in these four
un-divided rows, and its real form is
\[
 C_0W^2+A_0Y^2+
                 (D_0+\tfrac12)(W+vX)^2+\tfrac12\Gamma^2\succ0.
 \tag{WN42}
\]
The coefficient-extraction substitution \(X=R/v,Z=\Gamma-sR/v\)
cancels all denominators; the asserted identity itself contains
no \(v^{-1}\), and includes literal \(v=0\).

For both reciprocal angles \(s,b=o(1)\), the full polynomial
belongs to the square ideal of \((W,X,Y,sZ,bZ)\), with real form
\[
 C_0W^2+A_0Y^2+D_0(X-sZ)^2+
                    \tfrac12X^2+\tfrac12(Y-bZ)^2\succ0.
 \tag{WN43}
\]
These finite identities include \(s=0\) and \(b=0\) exactly.

It remains to check that the local state charts were exhaustive.
Normalize a maximal state coordinate in the testing field.
At \(u,v=o(1)\), a unit \(X\) or \(Y\) gives positive full
residue immediately. Otherwise one of \(W,Z\) is a unit.
The exact sign gauges give \(u,v\geq0\).
If \(WZ\leq0\), the full-core sign identity SE9 supplies
separated nonnegative terms; all event squares are positive.
If \(WZ>0\), choose \(W=1,Z\) bounded, or \(Z=1,W=o(1)\).
These are exactly the states covered above; vanishing angles
or states are included by the same polynomial identities
(or by the nonpositive-sign case).
Positive pure-field homogeneous normalization and real/imaginary
state splitting now restore the original full parent:
\[
 F_c(\hbox{any independent projective test and state})
       \ \in\ \Sigma(K_c\otimes_{\mathbb R}K_t)^2 .
 \tag{WN44}
\]
This proves precisely the weak-\(n\) C2 and all-\(\beta\) C1
weak-\(\alpha\) source domains in the source domain stated above.
It supplies the missing local leaves usually called WB, WU
and CW. It does not, by itself, assert the global source
partition or the finite-lift conclusion.



\section{\texorpdfstring{A source-exact projective pitchfork for C1}{A source-exact projective pitchfork for C1}}
\label{component:p2cdf6197677b}

\label{p2cdf6197677b:sec:new-c1-projective-pitchfork}

\paragraph{Theorem and precise source domain.}
Take a bounded physical canonical C1 source with $d=1$,
$\lambda=n=0$, $m,C,A,\alpha$ positive bounded units and
$D,\beta\geq0$ infinitesimal. Assume the original actual finite
noncommuting stationary zero contact, and retain every original
source gap and intersection. Then the entire parent has a finite
sum-of-squares certificate in $K_c\otimes_{\mathbb R}K_t$ for
every independent projective test and every real or complex state.
This includes literal $D=0$ or $\beta=0$ whenever allowed by the
source domain. The two cannot vanish together at an actual contact.

\subsubsection*{Opening-safe source entry}

Write the positive source contact as $(u_c,w_c)=(c,1/v_c)$,
and $V=w_c^2$. The retained positive-$d$ balance proof gives
$c$ a positive bounded unit and $V$ a positive infinitesimal.
Its second stationarity equation, in the notation of that proof, is
\[
 a_3=V^{-2}\{mDa_2/c^2+\beta ma_4+Da_5+\beta a_7c^2\},
 \tag{PP1}
\]
where all $a_i$ are positive bounded units. Therefore
\[
 D=V^2\delta,\quad\beta=V^2\eta,\qquad
 \delta,\eta\geq0\text{ bounded},\qquad
 \delta+\eta\text{ a positive unit}.
 \tag{PP2}
\]
No inverse of either opening is used.

Normalize the source state by $Z_c=1$. Set
\[
 p=W_c/w_c,\qquad y=Y_c/w_c,\qquad x=X_c/V.
 \tag{PP3}
\]
All three are bounded source elements. To prove this without a
shape hypothesis, restrict the reciprocal full parent PP6 to
$(W,Y,X)$ at $w=0$, and then take residues. Its positive hard form is
\[
 H_u(W,Y,X)=m(W+uY)^2+C(uW-Y)^2
       +(A+\alpha+1/2)(1+u^2)Y^2+(1+u^2)X^2/2.
 \tag{PP4}
\]
It is positive definite for every bounded real $u$. Scale the two
$W,Y$ kernel equations at the source by $w_c$, and the $X$
equation by $V$. In the unknowns $(p,y,x)$ the resulting bounded
matrix has a block triangular residue: its $p,y$ block is the
positive $W,Y$ block of PP4, and its $x$ pivot is $(1+c^2)/2$.
The terms $(\beta-D)/V=V(\eta-\delta)$ are bounded. All right
hand sides are bounded, proving PP3 by an ordinary source-field
matrix inverse with unit determinant.

Let $q_{01},q_{10},q_{21},h$ be the original Laurent coefficients.
Their first three residues are positive units because $m,C,A,\alpha$
are units. The actual zero and angular stationarity equations,
using PP2, give
\[
 \bar q_{01}=\bar q_{21}\bar c^{\,4},\qquad
 \bar h=-2\bar q_{21}\bar c^{\,2}.
 \tag{PP5}
\]
Every other nonconstant Laurent contribution apart from $q_{10}w^2$
vanishes in this residue calculation. These are equations wholly in
the coefficient field followed by its real residue map.

\subsubsection*{A complete local normal and divided-difference argument}

It remains locally to treat testing $w$ infinitesimal or zero and
$u$ with the same positive real residue as $c$. The full reciprocal
quarter-parent is
\[
\begin{split}
 F={}&m(W-wX+uY-uwZ)^2+CD_w(uW-Y)^2+AD_uD_wY^2\\
 &+(D+1/2)(uwW+uX-wY-Z)^2
       +(\alpha+1/2)D_u(Y-wZ)^2\\
 &+(\beta+1/2)D_w(X+uZ)^2-D_uD_wZ^2/2,
 \qquad D_u=1+u^2,\quad D_w=1+w^2.
\end{split}\tag{PP6}
\]
Use three unrestricted normals and two scalar conductors
\[
 N_W=W-pwZ,\quad N_Y=Y-ywZ,\quad N_X=X-xw^2Z,
 \qquad Q=w(u-c)Z,\quad L=(w^2-V)Z.
 \tag{PP7}
\]
The centers agree with the actual source contact; no equation is
imposed on an independent state. Write the full normal expansion
\[
 F=N^TH(u,w)N+2Z\sum_iN_i\ell_i(u,w)+E(u,w^2)Z^2,
 \quad N=(N_W,N_Y,N_X)^T.
 \tag{PP8}
\]
It is an exact expansion of every term of PP6. The matrix $H$ has
bounded entries and residue PP4 at $u=\bar c$. Parity and direct
inspection of PP6 give bounded polynomials $a_W,a_Y,a_X$ such that
\[
 \ell_W=w a_W(u,b),\qquad \ell_Y=w a_Y(u,b),\qquad
 \ell_X=u(\beta-D)+b a_X(u,b),\quad b=w^2.
 \tag{PP9}
\]
Actual state stationarity gives $a_W(c,V)=a_Y(c,V)=0$ and
$\ell_X(c,w_c)=0$. Finite monic division thus gives bounded
polynomials $A_i,B_i$ with
\[
 a_i(u,b)=(u-c)A_i(u,b)+(b-V)B_i(b),\qquad i=W,Y.
\]
Consequently the entire first two normal couplings are
$Z\ell_i=A_iQ+wB_iL$.

For the third normal put
$r(u)=u(\beta-D)/V+a_X(u,V)$. This is a bounded source
polynomial by PP2, and $r(c)=0$. Hence $r=(u-c)R(u)$ for
a bounded polynomial $R$. Division in $b$ gives another bounded
polynomial $B_X$ with
\[
 Z\ell_X=\{B_X(u,b)-(u-c)R(u)\}L+wR(u)Q.
 \tag{PP10}
\]
Here the identity $V(u-c)Z=wQ-(u-c)L$ is exact. Moreover
$B_X(c,0)$ has zero residue: the source equation
$u(\beta-D)+Va_X(u,V)=0$ at $u=c$, divided by $V$,
gives $a_X(c,V)$ infinitesimal, and $B_X(c,0)$ has the same
residue. Thus the third normal has no residue coupling to $L$
or $Q$. The first two have no residue coupling to $L$.

For completeness, the scalar coefficient obeys
\[
 E(u,0)=D+\beta u^2,\qquad
 E(c,V)=E_u(c,V)=E_b(c,V)=0.
 \tag{PP11}
\]
The first assertion follows by setting $w=0$ in the entire PP6
after the graph substitution. The other three follow from the
actual full source contact and its two angle derivatives; neither
normal term is discarded in making this inference.
Define $d(u)=\delta+\eta u^2$ and the bounded polynomial
\[
 J(u,b)=\{E(u,b)-d(u)(b-V)^2\}/b.
 \tag{PP12}
\]
This notation means exact monic division: PP11--PP2 make the
constant term zero. No testing $b$ is inverted. At $(c,V)$
the value and both first derivatives of $J$ vanish. A second
finite monic Taylor division gives bounded polynomials $J_{20}$,
$J_{11},J_{02}$ with
\[
 J(u,b)=(u-c)^2J_{20}(u,b)
       +(u-c)(b-V)J_{11}(b)+(b-V)^2J_{02}(b).
 \tag{PP13}
\]
Substitution yields the complete exact scalar identity
\[
 E(u,w^2)Z^2=J_{20}(u,w^2)Q^2+wJ_{11}(w^2)QL
                     +\{d(u)+w^2J_{02}(w^2)\}L^2.
 \tag{PP14}
\]
All divisions preceded tensor evaluation, with monic divisors and
bounded source coefficients. These formulas hold literally at $w=0$.

Equations PP8--PP14 represent the entire parent by a bounded
Gram on the five rows $N_W,N_Y,N_X,Q,L$. Its $L$ pivot in
residue is $\bar\delta+\bar\eta\bar c^2>0$, and its $L$
crosses have zero residue. It remains to check the four-row hard
and angular residue; the following exact determinant computation
provides that check without taking a mixed inverse.

\subsubsection*{Strictness of the angular residue}

At $D=\beta=0$, replace $Z$ by $T/w$ in the whole reciprocal
parent and take its polynomial value at $w=0$. This gives a real
four-state quadratic form $B_u(W,Y,X,T)$, with hard block PP4.
The original determinant identity gives
\[
 \det B_u=\frac{(1+u^2)^2}{16}
                  (q_{01}+hu^2+q_{21}u^4).
 \tag{PP15}
\]
Using PP5, its real Schur coefficient at $T$ is
\[
 g(u)=\frac{(1+u^2)^2q_{21}}{16\det H_u}
                       (u^2-c^2)^2.
 \tag{PP16}
\]
Here all entries and inverses are real residues; this calculation
proves positivity of a real matrix and does not factor a matrix in
the tensor ring. The contact centers PP3 give the unique real
kernel of $B_c$, so the four-row residue in PP8--PP14 is precisely
the quadratic Taylor coefficient in hard normals and the change
$u-c$. Its angular Schur pivot is
\[
 g''(c)/2=\frac{(1+c^2)^2q_{21}c^2}{4\det H_c}>0.
 \tag{PP17}
\]
The hard block is PP4 and is positive definite. Thus this full
four-row residue is positive definite, as is the independent $L$
pivot. The separated strict bounded-Gram lemma proves that PP6
is tensor SOS in the complete local chart. The rows may be
unbounded or linearly dependent at a literal endpoint; only the
coefficient Gram needs a strict real residue.

\subsubsection*{Explicit testing complement and projective endpoints}

If $u,v$ are bounded original testing coordinates, use the ordinary
state matrix except near the double zero corner. At finite nonzero
real residues, the residue Laurent polynomial is
\[
 q_{21}(u^2-c^2)^2/u^2+q_{10}/v^2>0,
 \tag{PP18}
\]
so its state matrix is positive definite by the hard $Z=0$ block.
At $u=0$, $v\ne0$ the real quarter-parent is
\[
 m(vW-X)^2+(C+A+1/2)(1+v^2)Y^2
                       +\alpha(vY-Z)^2+(1+v^2)X^2/2\succ0.
\]
At $v=0$, $u\ne0$ it is
$m(X+uZ)^2+(C+1/2)(uW-Y)^2+A(1+u^2)Y^2
+\alpha(1+u^2)Z^2+(X+uZ)^2/2$, which is likewise strict.
These are residue identities; the bounded full matrices therefore
give tensor SOS on their respective charts.

When both $u,v$ are infinitesimal or zero, use the five rows
$(uW,vW,X,Y,Z)$ in the entire original finite parent. The exact
polynomial substitution is bounded and its residue is
\[
 m(vW-X)^2+(C+1/2)(uW-Y)^2+AY^2+\alpha Z^2+X^2/2\succ0
 \tag{PP19}
\]
as a coefficient Gram on those five rows. This treats that corner
with no state normalization.

For reciprocal Alice and bounded Bob, the complete core identity
RC20 applies unchanged with all $m,D,\alpha,\beta$ event squares
retained whole. Its residue RC21 is strict using only $C,A>0$.
For two reciprocal testing angles, retain those positive squares
and use RC24 for the four rows $W,X,Y,wZ$. The residue
$CW^2+AY^2+X^2/2+(Y-wZ)^2/2$ is strictly positive.
These identities allow the reciprocal coordinate to be literally zero.

For Bob reciprocal and bounded Alice, retain the whole $D,\beta$
squares, put $T=wZ$, and use the full four-row expression
\[
\begin{split}
 m(W-wX+uY-uT)^2+CD_w(uW-Y)^2
          +AD_uD_wY^2+\alpha D_u(Y-T)^2+K/2,
\end{split}
\]
where $K$ is the full core RC22 with $\rho_*=uW-Y$.
Its residue is exactly $B_u$. By PP15--PP16 it is strictly
positive whenever $u^2$ has residue different from $\bar c^2$,
including $u$ infinitesimal. The remaining matched residues are
PP7--PP17 after the exact Alice sign gauge if necessary.

Both signs of either angle are allowed by the exact simultaneous
state sign gauges of the original parent. Every normalization is
a positive square in a single source or testing field. Complex
states split into real and imaginary parts. The finite and
reciprocal homogeneous identities include all literal axes and
projective endpoints. This completes the all-tests theorem on the
stated positive-$d$ C1 projective source domain.



\section{\texorpdfstring{A direct certificate for the positive-$d$ C2 reciprocal corner}{A direct certificate for the positive-d C2 reciprocal corner}}
\label{component:p10686597065a}

\label{p10686597065a:sec:new-rc-direct}

\paragraph{The source domain.}
Normalize the original C2 source by $d=1$ and $\lambda=\beta=0$.
Assume that all source coefficients are bounded, that $C,A,\alpha$
are positive bounded units, that $m,D$ are positive infinitesimals,
and that $n\geq0$ is bounded. Retain the original physicality,
finite noncommuting stationary zero contact and positive native gaps. This section proves
the complete parent is tensor SOS at every independent projective
test and every state. The contact is used only inside the source
field. No mixed determinant is cancelled.

Equation~\eqref{eq:reciprocal-source-balances}, with $\beta=0$,
gives, for the reciprocal Bob contact coordinate $w_c=1/v_c$,
\[
 U=u_c^2\in\operatorname{Inf}_{>0},\qquad
 V=w_c^2\in\operatorname{Inf}_{>0},\qquad
 m\asymp U^2,\qquad D\asymp V^2.
 \tag{RC1}
\]
The positive-unit ratios follow from the final formula in
Section~\ref{lem:reciprocal-source-balance}.

\subsubsection*{The actual source has a bounded shape chart}

Choose the native positive contact chart with $Z_c=1$, and set
\[
 S=(u_c/w_c)W_c,\qquad y=Y_c/w_c,\qquad
 R=-X_c/(u_cV).
 \tag{RC2}
\]
Then $S,y,S-y,1-y,R$ are positive bounded units. Here is an entry
proof that does not assume those facts as extra source hypotheses.
In the reciprocal parent RC6 below, take the ordinary source-field
Schur complement of the $X,Y$ block. That block has residue
$\operatorname{diag}(n+1/2,H)$, where $H=C+A+\alpha+1/2$.
Its $W,W$ Schur entry divided by $U$ has residue
$a=C-C^2/H>0$, while its $W,Z$ entry divided by $u_cw_c$
has residue $-b=-(1/2+C\alpha/H)<0$. The balance
$m/U\in\operatorname{Inf}$ and $D/V\in\operatorname{Inf}$
justifies these divisions in the source field. The zero-state
equation therefore gives $S\sim b/a$, a positive bounded unit.

The two hard state equations, scaled in the source field by
$u_cV$ and $w_c$, are a bounded linear system in $R,y$. Its
residue is triangular with positive diagonal $n+1/2,H$.
Consequently both quantities are bounded. The second equation gives
$Hy=CS+\alpha$ in residue, so $y$ is a positive unit. The
$W$ equation divided by $u_cw_c$ gives
$C(S-y)=1/2$ in residue. The $Z$ equation divided by $V$
gives $\alpha(1-y)=S/2$ in residue: its apparently constant
$U$ terms cancel exactly before the division. Finally the first
hard equation gives $(n+1/2)R=(1+S-y)/2$ in residue. These
relations prove all five claimed positive-unit bounds. The exact
native gaps already give $0<y<1$ and $S-y>0$ before taking
residues.

The unchanged native source formulas are converted by the following
dictionary. Their letters are used only in this display:
\[
\begin{aligned}
 z=U,\quad a^2=1/(UV),\quad b=S/U,\quad t=(S-y)V,
 &\quad r=VR,\quad h=y/U,\\
 g=S/U+y-1+VR,&\qquad j=1+UVR-V(S-y).
\end{aligned}
 \tag{RC3}
\]
Substitution into the six exact native coefficients has denominators
with positive-unit leading parts after the displayed powers of
$U,V$ are removed. Thus $\mu=m/U^2$ and $\delta=D/V^2$
are bounded units, and the source residues are
\[
\begin{aligned}
 C_0&=\frac1{2(S-y)},& A_0&=\frac{1+S-y}{2y},&
 \alpha_0&=\frac{S}{2(1-y)},\\
 n_0&=\frac{1+S-y-R}{2R},&
 \mu_0&=\frac{1+S-y}{2S^2},&
 \delta_0&=\frac{S(1+S-y)}2.
\end{aligned}\tag{RC4}
\]
In residue formulas such as RC4 the shape letters mean their real
residues. In particular
\[
 C_0\{\alpha_0(A_0-1/2)-1/4\}
       =(A_0+\alpha_0+1/2)/4.
 \tag{RC5}
\]
These are source statements only. They do not identify any source
infinitesimal with an independent testing infinitesimal.

\subsubsection*{The complete local parent and both unrestricted normals}

For $u,w$ infinitesimal or zero in the independent testing field,
put $D_u=1+u^2$, $D_w=1+w^2$. The full quarter-parent after
the pure testing reciprocal Bob normalization is
\[
\begin{split}
 F={}&m(W-wX+uY-uwZ)^2+nD_uD_wX^2
       +CD_w(uW-Y)^2+AD_uD_wY^2\\
 &+(D+1/2)(uwW+uX-wY-Z)^2
       +(\alpha+1/2)D_u(Y-wZ)^2\\
 &+\tfrac12D_w(X+uZ)^2-\tfrac12D_uD_wZ^2.
\end{split}\tag{RC6}
\]
All four state coordinates are arbitrary. Introduce tensor rows
\[
 \begin{gathered}
 N_x=X+uRw^2Z,\qquad N_y=Y-ywZ,\\
 \rho=uW-SwZ,\qquad k=UW-SuwZ,\qquad \ell=(w^2-V)Z.
 \end{gathered}
 \tag{RC7}
\]
The testing-dependent $w^2$ in $N_x$ is essential; it agrees with
the source contact and removes the unbounded coefficient produced
by a contact-fixed center. The exact conductor identities are
\[
 U\rho-uk=Sw(u^2-U)Z,\qquad
 u\rho-k=(u^2-U)W.
 \tag{RC8}
\]

Expanding the entire RC6 gives
\[
 F=F_0+H_xN_x^2+H_yN_y^2+2H_{xy}N_xN_y
                         +2N_xL_x+2N_yL_y,
 \tag{RC9}
\]
where $F_0$ is only the constant coefficient of this unrestricted
normal expansion. The hard coefficients are exactly
\[
\begin{split}
 H_x&=mw^2+nD_uD_w+(D+1/2)u^2+D_w/2,\\
 H_y&=mu^2+CD_w+AD_uD_w+(D+1/2)w^2+(\alpha+1/2)D_u,\\
 H_{xy}&=-(m+D+1/2)uw.
\end{split}\tag{RC10}
\]
Set $J_x=n(1+w^2)+D+1/2$ and $K_x=n(1+U)+1/2+m$.
The complete first normal is
\[
 L_x=(D+1/2)uw\rho-(m/U)wk
 +u\{D/V-Rw^2K_x\}\ell
 -(R/S)uwJ_x(U\rho-uk).
 \tag{RC11}
\]
Every coefficient on $\rho,k,\ell$ is bounded infinitesimal.
For the other normal define
\[
 \begin{aligned}
 a_{10}&=Ay+(m+\alpha+1/2)(y-1),\\
 a_{01}&=-(C+D+1/2)S+(C+A+D+1/2)y,\\
 a_{11}&=R(m+D+1/2)+Ay.
 \end{aligned}
\]
Then the exact expression is
\[
\begin{split}
 L_y={}&\{m-C-(C+D+1/2)w^2+(U/S)(a_{10}+a_{11}w^2)\}\rho\\
 &-(u/S)(a_{10}+a_{11}w^2)k+w(a_{01}+a_{11}U)\ell.
\end{split}\tag{RC12}
\]
Its only nonzero residue coupling is $-C_0\rho$. RC11--RC12
follow by full expansion and the two actual source state equations;
they account for every term linear in the two normals in RC6.

\subsubsection*{A scalar certificate by finite monic divisions}

Put
\[
 c=C+(C+D+1/2)w^2,\quad b_2=(D+1/2)R,
 \quad b_1=(C+D+1/2)(y-S)-mR,\quad \gamma=b_1+b_2U.
\]
Define
\[
 F_{\rm base}=(c-2Ub_2w^2/S)\rho^2+\mu k^2
       +(2ub_2w^2/S)\rho k-2w\gamma\rho\ell.
 \tag{RC13}
\]
The source $W$ state equation removes the remaining $WZ$
coefficient. Direct expansion then gives
$F_0-F_{\rm base}=E(u^2,w^2)Z^2$, with $E(a,b)$ a
bounded-coefficient polynomial of degree at most two in $a$.
Its exact properties are
\[
\begin{gathered}
 E(a,0)=D,\qquad E(U,V)=E_a(U,V)=E_b(U,V)=0,\\
 E(a,b)=E_0(b)+(a-U)E_1(b)+(a-U)^2E_2(b),\\
 E_2(b)=R^2b^2\{n(1+b)+D+1/2\}.
\end{gathered}\tag{RC14}
\]
The three vanishing jets follow from the actual full contact:
every row in RC7 vanishes there, the complete scalar and both
angle derivatives vanish, and RC13 has zero first derivatives
there. They are also checked directly from the native source
coefficients. Since $E_1(0)=E_1(V)=0$ and $E_0$ has a double
zero at $V$, finite monic polynomial division gives bounded
polynomials $J_0,J_1$ such that
\[
 E_0(b)=(b-V)^2J_0(b),\qquad
 E_1(b)=b(b-V)J_1(b).
 \tag{RC15}
\]
These divisions take place in the source polynomial algebra before
tensor evaluation. They introduce no inverse of $V$ or of a mixed
quantity. In particular $V^2J_0(0)=D$, and the other coefficients
of $J_0$ remain bounded; hence $J_0(w^2)$ has residue $\delta_0$.

RC8 now gives the exact scalar identity
\[
\begin{split}
 F_0={}&F_{\rm base}
 +\frac{R^2w^2\{n(1+w^2)+D+1/2\}}{S^2}(U\rho-uk)^2\\
 &+\frac wS J_1(w^2)(U\rho-uk)\ell+J_0(w^2)\ell^2.
\end{split}\tag{RC16}
\]
This polynomial formula also holds at $u=0$ or $w=0$; neither
is inverted. Combining RC9--RC16 represents the entire parent
by a bounded Gram on $(N_x,N_y,\rho,k,\ell)$ with real residue
\[
 (n_0+1/2)N_x^2+H_0N_y^2-2C_0N_y\rho+C_0\rho^2
                              +\mu_0k^2+\delta_0\ell^2,
 \qquad H_0=C_0+A_0+\alpha_0+1/2.
 \tag{RC17}
\]
This is positive definite: the $N_y,\rho$ determinant is
$C_0(A_0+\alpha_0+1/2)>0$, and the other three pivots are
strictly positive. The separated strict bounded-Gram lemma
therefore gives a finite tensor SOS of the complete RC6. It
requires bounded matrix entries, not bounded rows or states.

\subsubsection*{All other testing charts}

The remaining measurement charts follow from three matrix
identities. In each one we retain the indicated event squares and
prove positivity of the remaining bounded Gram matrix.

\paragraph{Both original testing coordinates bounded.}
Put $p=uW$. Retain the whole $m,n,D$ squares and use the
remaining Gram on $(p,Y,X,Z)$:
\[
\begin{split}
 C D_v(p-Y)^2+A D_uD_vY^2
 &+\tfrac12(p+uvX-Y-vZ)^2+(\alpha+1/2)D_u(vY-Z)^2\\
 &+\tfrac12D_v(X+uZ)^2-\tfrac12D_uD_vZ^2.
\end{split}\tag{RC18}
\]
The $Z=0$ hard block is positive definite by its displayed
$C,A,X$ squares. After RC4 the determinant of the real residue is
\[
 \frac{(S+u^2v^2)(1+u^2)^2(1+v^2)^2(1+S-y)^2}
      {16y(S-y)(1-y)}>0.
 \tag{RC19}
\]
Here also $u,v$ mean real testing residues. Thus RC18 is a strict
bounded Gram, even when $u$ or $v$ has zero residue.

\paragraph{Alice reciprocal, Bob bounded.}
Let $s=1/u$ be infinitesimal or zero, set $D_s=1+s^2$, and use
rows $(W,Y,V_x,R_x)=(W,Y,vX,sX+Z)$. Retain the entire
$m,n,D,\alpha$ squares. The remaining expression is
$CD_v(W-sY)^2+AD_sD_vY^2+K/2$, where exactly
\[
\begin{split}
 K={}&\{W-sY+D_sV_x-svR_x\}^2+D_vR_x^2\\
 &+D_s\{v^2Y^2-2Y(vR_x-sV_x)-(vR_x-sV_x)^2\}.
\end{split}\tag{RC20}
\]
Its residue is
\[
 C_0D_vW^2+A_0D_vY^2
          +\tfrac12(W+V_x)^2+\tfrac12(vY-R_x)^2\succ0.
 \tag{RC21}
\]
This single formula includes all bounded $v$, including literal
zero and arbitrary competing small $s,v$. There is no $s/v$
coefficient.

\paragraph{Bob reciprocal, Alice a bounded unit.}
Let $w$ be infinitesimal or zero, and $u$ a bounded unit. Use
$\rho_*=uW-Y$, $T=wZ$, $S_D=\sqrt D\,Z$, and the five rows
$(\rho_*,Y,X,T,S_D)$. Retain $m,n$ whole. The $D$ square is
exactly $\{\sqrt D(uX+w\rho_*)-S_D\}^2$. The other terms are
\[
 CD_w\rho_*^2+AD_uD_wY^2+\alpha D_u(Y-T)^2+K/2,
\]
where the complete core is
\[
 K=(D_u+w^2)X^2+2uwX T+w^2\rho_*^2+2uw\rho_*X
                           -2\rho_*T-T^2+D_u(Y-T)^2.
 \tag{RC22}
\]
The $X,S_D$ residue pivots are $D_u/2,1$. On
$(\rho_*,Y,T)$ the residue matrix is
\[
 \begin{pmatrix}
 C_0&0&-1/2\\
 0&G D_u&-(\alpha_0+1/2)D_u\\
 -1/2&-(\alpha_0+1/2)D_u&\alpha_0D_u+u^2/2
 \end{pmatrix},\qquad G=A_0+\alpha_0+1/2.
 \tag{RC23}
\]
Its first two leading minors are $C_0,C_0GD_u$. By RC5 its
determinant is $GD_uu^2(1+2C_0)/4>0$. This proves strictness.
The alternative $u$ infinitesimal or zero is precisely RC7--RC17.

\paragraph{Both testing coordinates reciprocal.}
Let $s,w$ be infinitesimal or zero. Set $T=wZ$ and
$S_D=\sqrt D\,sZ$. Use the five rows $W,X,Y,T,S_D$, and retain the
entire $m,n,\alpha$ squares. The $D$ square is
$\{\sqrt D(wW+X-swY)-S_D\}^2$. The other terms are
$CD_w(W-sY)^2+AD_sD_wY^2+K/2$, where exactly
\[
\begin{split}
 K={}&(wW+X-swY)^2+D_s(Y-T)^2+s^2D_wX^2\\
 &-2sWT+2s^2YT+2swXT-s^2T^2.
\end{split}\tag{RC24}
\]
The full five-row residue is
\[
 C_0W^2+A_0Y^2+X^2/2+(Y-T)^2/2+S_D^2\succ0.
 \tag{RC25}
\]
This gives the final strict bounded Gram.

Every ordered-field testing coordinate is bounded or has
infinitesimal reciprocal. RC18--RC25 and the local proof therefore
cover both projective lines. The identities are polynomial in the
chosen homogeneous coordinates, so literal infinity is included by
$s=0$ or $w=0$. Both signs are allowed throughout; no state ratio
or state normalization is used. Positive pure testing square
normalizations restore the original homogeneous parent, and the
original source scale is restored by a pure source square root.
Real and imaginary parts of a complex state are handled separately.


